\documentclass[12pt]{article}
\usepackage[utf8]{inputenc}
\renewcommand{\baselinestretch}{1.5} 

\newcommand*{\SuperScriptSameStyle}[1]{%
  \ensuremath{%
    \mathchoice
      {{}^{\displaystyle #1}}%
      {{}^{\textstyle #1}}%
      {{}^{\scriptstyle #1}}%
      {{}^{\scriptscriptstyle #1}}%
  }%
}

\newcommand{\oneS}{\ensuremath{{}^{\textstyle *}}}
\newcommand{\twoS}{\ensuremath{{}^{\textstyle **}}}
\newcommand{\threeS}{\ensuremath{{}^{\textstyle **}\oneS}}

\usepackage{amsmath} % for equations,  math functions and operators
\usepackage{amsfonts} % extended fonts
\usepackage{eurosym} % Euro symbol
\usepackage[hmargin=2.5cm,vmargin=3.5cm]{geometry}
\usepackage{graphicx} % \includegraphics preferable as pdf or png
\graphicspath{{images/}{figures/}{figures/layout/}{figures/analysis/}}
%\usepackage{subfigure} % subfigures in subfigure environment
\usepackage[font+=small,labelfont=bf, skip=0pt]{caption,subcaption}
\usepackage{float}
%\usepackage{gensymb}
\usepackage{newunicodechar}
\newunicodechar{₂}{$_2$}
\newunicodechar{≈}{\approx}
\usepackage{newunicodechar}
\newunicodechar{ }{~}
%[labelfont=bf]
%\captionsetup{font=footnotesize}
\usepackage{multirow}
\usepackage{float}
\usepackage{graphicx}
\usepackage{color} % color text
\usepackage{setspace} % adjust line spacing with \onehalfspacing, \doublespacing and \singlespacing
\usepackage{footmisc} % adjust footnotes
\usepackage{natbib} % advanced bibliography management
\usepackage{pdflscape} % used for singular landscape pages
\usepackage{array} % visualizes matrixes 
\usepackage[debugshow]{tabulary}
\usepackage[debugshow]{tabularx}
\usepackage{amssymb}
\usepackage{ragged2e}
\usepackage[para,online,flushleft]{threeparttable}
\usepackage[etex=true,export]{adjustbox}
\usepackage{tablefootnote}
\usepackage{tablefootnote}
\usepackage[colorlinks=true,allcolors=blue]{hyperref} % set hyperlinks in the document and to the outside (webpages, mail-adrdresses)
\usepackage{ntheorem} % replacement for  amsthm to use \theoremstyle{break}
\usepackage{tikz}
\usepackage{lscape}
\usepackage{comment}
\usepackage{mathtools}
\usepackage{dsfont}
\usepackage{multimedia}
\usepackage{adjustbox}
%\usetikzlibrary{shapes,snakes}
%\usetikzlibrary{arrows.meta}
\usepackage{varwidth}
\usepackage{soul}
\usetikzlibrary{positioning}
%\usetikzlibrary{shapes.multipart}
\newenvironment{changemargin}[2]{%
	\begin{list}{}{%
			\setlength{\topsep}{0pt}%
			\setlength{\leftmargin}{#1}%
			\setlength{\rightmargin}{#2}%
			\setlength{\listparindent}{\parindent}%
			\setlength{\itemindent}{\parindent}%
			\setlength{\parsep}{\parskip}%
		}%
		\item[]}{\end{list}}

\usepackage{tikz}
\usetikzlibrary{trees}

\tikzstyle{level 1}=[level distance=3.5cm, sibling distance=4.0cm]
\tikzstyle{level 2}=[level distance=3.5cm, sibling distance=2cm]

\tikzstyle{bag} = [align=center]

\usepackage{mathtools}
\usepackage{booktabs}
\usepackage{longtable}
\usepackage[hmargin=2.5cm,vmargin=3.5cm]{geometry}
\DeclarePairedDelimiter\ceil{\lceil}{\rceil}
\DeclarePairedDelimiter\floor{\lfloor}{\rfloor}

% For tick boxes in the appendix: 
\usepackage{enumerate}
\usepackage{amssymb}
\usepackage[inline]{enumitem}
\newlist{todolist}{itemize}{2}
\setlist[todolist]{label=$\square$}
\usepackage{pifont}
\newcommand{\cmark}{\ding{51}}%
\newcommand{\xmark}{\ding{55}}%
\newcommand{\done}{\rlap{$\square$}{\raisebox{2pt}{\large\hspace{1pt}\cmark}}%
\hspace{-2.5pt}}
\newcommand{\wontfix}{\rlap{$\square$}{\large\hspace{1pt}\xmark}}

%Trade mark sign
\usepackage{textcomp}
\usepackage{diagbox}

%\doublespacing
\onehalfspacing
\theoremstyle{break} % options: break (default), definition, plian, remark
% Here you can definition your own theorem environments. DON'T USE A HIERACHY!
\newtheorem{definition}{Definition}
\newtheorem{assumption}{Assumption}

\newtheorem{theorem}{Theorem}
\newtheorem{corollary}{Corollary}
\newtheorem{lemma}{Lemma}
\newtheorem{proposition}{Proposition}

% Always make a linespace in front of every proof.
\newenvironment{proof}[1][Proof]{\noindent\textbf{#1.}}{\ \rule{0.5em}{0.5em}}

% Use \C{width}, \L{widht} and \R{width} in tabular environment to make tables with fixed column width. C, L, and R stand for center, align left and align right
\newcolumntype{L}[1]{>{\raggedright\let\newline\\arraybackslash\hspace{0pt}}m{#1}}
\newcolumntype{C}[1]{>{\centering\let\newline\\arraybackslash\hspace{0pt}}m{#1}}
\newcolumntype{R}[1]{>{\raggedleft\let\newline\\arraybackslash\hspace{0pt}}m{#1}}

% adjust your margins here
\geometry{
	left=2.54cm,
	right=2.54cm,
	top=2.54cm,
	bottom=2.54cm
}

%\usepackage[font=small,labelfont=bf,skip=0pt]{caption}
\usepackage{dcolumn}
\usepackage[toc,page]{appendix}
\renewcommand\appendixpagename{Online Appendix}
\usepackage[noabbrev,nameinlink,capitalise]{cleveref}

\providecommand{\keywords}[1]{\textbf{Keywords:} #1}
\providecommand{\jelcodes}[1]{\textbf{JEL codes:} #1}
\crefname{appsec}{Section}{Appendices}

% --- ETH title page helpers ---
\usepackage{titling}   % lets us access \thetitle, \theauthor, etc.
\usepackage{graphicx}
\usepackage{setspace}

% Thesis metadata (edit these)
\newcommand{\thesistype}{Master Thesis}
\newcommand{\studentname}{Gian-Luca Kaufmann}
\newcommand{\studentemail}{gikaufmann@ethz.ch}
\newcommand{\studentnumber}{20-913-166}
\newcommand{\thesistitle}{Predicting Protected Area Designations and Assessing Transition Risk under 30x30}
\newcommand{\thesissubtitle}{}        
\newcommand{\chair}{Chair of Energy and Climate Economics}
\newcommand{\department}{Department of Management, Technology, and Economics, ETH Zürich}
\newcommand{\supervisorblock}{Prof.\ Dr.\ Lint Barrage, D-MTEC ETH Zurich \\
Dr.\ Patricio Velasquez, D-MTEC ETH Zurich \\
Dr.\ Chiara Colesanti Senni, University of Zurich\\
Elena\ Almeida, CETEx LSE}
\newcommand{\submissiondate}{January 2026} % edit

\begin{document}
\begin{titlepage}
\thispagestyle{empty}

% Top logos
\vspace*{-1.5cm}
\noindent
\begin{minipage}[t]{0.49\textwidth}
  
\includegraphics[height=1.5cm]{figures/layout/eth_logo_black.pdf}
\end{minipage}
\begin{minipage}[t]{0.49\textwidth}
  \raggedleft
  
\includegraphics[height=1cm]{figures/layout/LOGO_MTEC.png}
\end{minipage}

% Main title block
\vspace*{2cm}
\begin{center}
{\LARGE\bfseries \thesistitle \par}
\vspace{0.4cm}
{\Large\bfseries \thesissubtitle \par}
\vspace{0.8cm}
{\large \thesistype \par}
\end{center}

% Author block
\vspace*{2.0cm}
\begin{center}
{\large \studentname \par}
\vspace{0.25cm}
{\small \studentemail \par}
\vspace{0.25cm}
{\small \studentnumber \par}
\vspace{1.2cm}
{\small \submissiondate \par}
\end{center}

% Bottom institute + supervisors
\vfill
\begin{center}
{\small \chair \par}
{\small \department \par}
\vspace{1.4cm}
{\small \textbf{Supervisors:} \par}
{\small \supervisorblock \par}
\end{center}

\end{titlepage}

\singlespacing
\begin{abstract}
\noindent This thesis investigates the establishment of protected areas as a critical component of global environmental conservation efforts. The 30x30 target, which aims to protect 30\% of the planet's land and oceans by 2030, represents a significant step towards safeguarding biodiversity. However, uncertainty surrounding the specific locations of new protected areas generates substantial transition risk for various stakeholders, including policymakers, businesses and investors. Machine learning models can predict where new protected areas might be designated and help mitigate these risks. These models can analyse large and complex datasets to identify patterns and predict outcomes with a high degree of accuracy. In this thesis, such models are used to forecast potential locations for new protected areas in the context of achieving the 30x30 target in South America. By integrating a wide range of variables in a 1x1km resolution, including current land use, environmental value, existing protected regions, and regional economic activities, the models provide insights into likely future designations. This predictive capability enables stakeholders to anticipate and adapt to changes that could affect their operations and investments. A key factor considered in the models is the location of highly profitable activities, reflecting the assumption that protected areas are less likely to be established in regions with economically significant land uses such as mining, agriculture, or urban development. By incorporating economic data and land-use patterns, predictions can be refined and transition risk for businesses operating in these areas, as well as for prudential supervisors assessing such risks, can be reduced. This approach allows companies to plan more effectively and mitigate potential disruptions, while also supporting governments in planning protected area establishments and managing associated transition risks.
\end{abstract}
\onehalfspacing

\newpage
\tableofcontents
\newpage

\setstretch{1.5}
\newpage
\section{Introduction}

The global expansion of protected areas (PAs) is a central pillar of contemporary biodiversity policy. Under the “30×30” target, governments have committed to protecting 30\% of the world’s land and oceans by 2030, implying a substantial increase relative to current coverage (Convention on Biological Diversity, 2022). While this expansion is widely regarded as necessary to halt biodiversity loss, it also entails significant economic and regulatory consequences. The designation of new protected areas typically introduces binding land-use restrictions, affecting activities such as agriculture, mining, infrastructure development, and urban expansion (Schuster et al., 2023). As a result, uncertainty about where new protected areas will be established generates material \emph{transition risk} for businesses, investors, and public authorities, as future conservation decisions may strand assets or alter the viability of existing land uses.

Recent economic research highlights that protected area policies involve real opportunity costs in the form of foregone economic activity and constrained resource extraction (Reynaert et al., 2024). These costs are not evenly distributed across space, but depend critically on the specific locations chosen for protection. From a risk-management perspective, the key challenge is therefore not only the scale of protected area expansion, but its spatial allocation (Robinson et al., 2024). Anticipating where new PAs are likely to be designated would allow affected stakeholders to adapt investment strategies, land-use planning, and regulatory frameworks in advance, reducing the likelihood of abrupt economic disruptions. Reflecting this concern, nature-related financial risk frameworks explicitly identify the creation of protection efforts as a source of transition risk that can affect business operations and supply chains (Govind \& Webb, 2025).

Against this background, recent research has begun to explore whether machine learning (ML) methods can be used to predict the spatial expansion of protected areas. For example, Mouillot et al. (2024) models the “socio-environmental niche” of protected areas using a set of spatial covariates to identify locations with favourable conditions for PA designation. By training classification models on a global grid of protected and unprotected cells derived from the World Database on Protected Areas, ML models can successfully identify areas with a high likelihood of future protection, conditional on observed historical patterns. Such approaches provide an important proof of concept, showing that PA establishment is not random in space and can, to some extent, be anticipated.

However, existing ML-based approaches to predicting protected area establishment remain limited in their ability to inform transition risk assessments. Current models are typically designed to characterise the ecological or institutional conditions associated with protected areas, rather than to explicitly capture the economic and regulatory dimensions that are most relevant for transition risk. As a result, there is currently no modelling framework in the academic literature that directly predicts protected area establishment using a feature set and evaluation strategy tailored to transition risk considerations. This limitation leaves a critical gap between conservation-oriented predictive modelling and the practical needs of stakeholders exposed to policy-driven land-use change.

This thesis addresses this gap by developing a machine learning framework to predict the establishment of new protected areas in the context of global conservation initiatives such as the 30×30 target, with a specific focus on transition risk relevance. The study aims, first, to construct ML models that forecast future PA designations based on a comprehensive set of spatial predictors capturing environmental, land-use, and economic characteristics. Second, it evaluates how accurately these models can identify high-risk locations where future protection is most likely, thereby providing actionable information for businesses, investors, and policymakers facing regulatory and economic uncertainty associated with conservation expansion.
\newpage 

\section{Literature Review}\label{literature-review}

\subsection{Global Biodiversity Decline and the Drivers of Ecological Pressure} \label{biodiversity-decline}

Human activities have profoundly altered natural systems, leading to a rapid decline in the capacity of ecosystems to provide essential services, commonly referred to as ecosystem services (Díaz et al., 2019). Biodiversity plays a central role in sustaining these services, yet global biodiversity loss is occurring at an unprecedented rate in human history (Brondizio et al., 2019). The erosion of biological diversity weakens ecosystem functioning and reduces the ability of natural systems to buffer environmental change and natural disturbances (Hooper et al., 2005). Beyond ecological consequences, biodiversity loss poses direct risks to human well-being, including reduced food security and increased pollution of soil and water resources (Xu et al., 2021).

The drivers of biodiversity decline are overwhelmingly anthropogenic. Human activities have introduced multiple, interacting pressures on natural ecosystems, including land-use change, overexploitation of natural resources, pollution, the spread of invasive species, and climate change (Millennium Ecosystem Assessment, 2005). Among these drivers, climate change has emerged as a particularly powerful force accelerating biodiversity loss. It is already catalysing large-scale ecological disruptions, such as the melting of the Greenland ice sheet and the bleaching of coral reef systems including the Great Barrier Reef (Nicholls and Cazenave, 2010; Hughes et al., 2017). Freshwater ecosystems, in particular, are experiencing widespread biodiversity pressure as a combined result of human intervention and climate change (Jeremias et al., 2018).

Climate change also affects human systems, reinforcing the interconnected nature of ecological and socio-economic risks. Changes in temperature and precipitation patterns threaten global food security and the availability of drinking water (Wheeler and von Braun, 2013). These dynamics highlight the growing trade-offs and synergies between climate change mitigation, biodiversity conservation, and land-use decisions, underlining the importance of effective land management at the global scale (Prudhomme et al., 2020). While climate change is expected to cause further biodiversity loss, the conservation of biodiversity itself can act as a climate mitigation mechanism. Protecting ecosystems contributes directly to climate regulation and can reduce rising temperatures (Shin et al., 2022). Consistent with this view, Girardin et al. (2021) show that nature-based solutions, including ecosystem protection, can lead to significant global cooling effects.

In addition to climate change, land-use change constitutes a major and direct threat to biodiversity. Socio-economic development processes such as urbanisation and agricultural expansion strongly influence land-use and land-cover patterns, leading to habitat fragmentation, soil degradation, and deforestation (Schippers et al., 2015). Empirical studies indicate that conventional land-use change is often associated with land-use intensification, which tends to increase agricultural production at the expense of biodiversity (Beckmann et al., 2019). Urbanisation similarly exerts negative effects on species diversity, particularly at high levels of urban development (McKinney, 2008). While some aspects of land-use change remain contested, for example, the ecological effects of habitat fragmentation continue to be debated (Valente et al., 2023), the impact of habitat loss itself is well established and consistently associated with biodiversity decline (Fahrig, 2003).

\subsection{Protected Areas: Purpose, History, and Global Trends} \label{pa-purpose-history}

In response to the accelerating loss of global biodiversity, the establishment of protected areas has become a central policy instrument worldwide (Juffe-Bignoli et al., 2014). Protected areas are intended to safeguard ecosystems and biodiversity by placing legal or regulatory constraints on land and resource use. The International Union for Conservation of Nature (IUCN) defines protected areas as “clearly defined geographical space[s], recognised, dedicated and managed, through legal or other effective means, to achieve the long-term conservation of nature with associated ecosystem services and cultural values” (Dudley and Stolton, 2008). In practice, this overarching definition is complemented by national policies and legal frameworks, which determine how protected areas are designated and governed within individual countries (UNEP-WCMC and IUCN, 2024).

The effectiveness of protected areas in conserving biodiversity and ecosystem services depends critically on the political, institutional, and socio-economic contexts in which they are established and managed. In many regions, underfunding, inadequate management capacity, corruption, and armed conflict undermine the ability of protected areas to achieve their conservation objectives (Watson et al., 2014). Moreover, protected areas do not function in isolation. They are embedded within broader landscapes and seascapes and interact continuously with surrounding areas. As a result, protected areas can act as sources of ecological benefits to neighbouring regions while simultaneously being exposed to external environmental pressures (Palomo et al., 2013).

Ecological processes frequently transcend protected area boundaries. Many species move across protected area borders on daily, seasonal, or nomadic timescales, which exposes them to threats such as exploitation and land-use change occurring outside protected zones (Gatson et al., 2008). Consequently, some of the most significant threats to protected areas originate beyond their formal boundaries and are driven by broader economic, socio-political, and land-use dynamics. These external pressures can degrade ecosystems within protected areas themselves. At the same time, ecosystem services generated within protected areas often flow outward, benefiting surrounding land uses and populations, with nearby urban areas in some cases depending directly on these services (Palomo et al., 2013).

While much of the literature focuses on the management and effectiveness of protected areas, this thesis is concerned primarily with their establishment, and in particular with the spatial location of new protected areas. The distributional consequences of protected area establishment for local communities remain insufficiently understood. Evidence suggests that the costs and benefits of protected areas are not evenly shared. For example, Ma et al. (2018) find that biodiversity conservation in protected reserves imposes significantly higher direct economic costs on households located within reserve boundaries compared to those outside, while direct economic benefits do not differ significantly. In addition, indirect costs associated with protected areas can be substantial. Yang et al. (2018) show that increased human–wildlife conflict linked to afforestation can impose significant indirect burdens on affected communities.

Overall, the literature indicates that protected areas have wide-ranging effects that extend beyond biodiversity conservation alone. By providing habitats for biodiversity, protected areas directly contribute to the achievement of the United Nations Sustainable Development Goals (SDGs), particularly Goal 14 (Life Below Water) and Goal 15 (Life on Land). They also influence Goal 13 (Climate Action) indirectly through the cooling effects of vegetation. Furthermore, existing research highlights links between protected areas and other development objectives, including Goal 6 (Clean Water and Sanitation) and Goal 2 (Zero Hunger), through their roles in maintaining water cycles and supporting food security.
\newline \newline
Protected areas have existed in various forms for millennia. Long before the emergence of modern conservation policies, many Indigenous and local communities established protected spaces through practices such as sacred or guarded sites, communal resource-use areas, and hunting grounds reserved for ruling classes. These early forms of protection reflect longstanding societal efforts to regulate access to natural resources and preserve areas of ecological, cultural, or symbolic importance. (Chape et al., 2005)

Protected areas in their modern institutional form emerged primarily during the nineteenth century in regions such as Europe, North America, Australia, New Zealand, and South Africa. During this period, the motivations for establishing protected areas varied considerably across regions. In Europe, protected areas were often created to preserve cultural landscapes, whereas in Africa the establishment of game parks was driven largely by wildlife conservation objectives. In North America, by contrast, the protection of dramatic natural scenery played a central role in early conservation efforts. Over the course of the twentieth century, the idea of protected areas spread globally, with an increasing number of countries adopting some form of area-based conservation. (Phillips, 2007)

As environmental degradation intensified in the latter half of the twentieth century, efforts to coordinate conservation policies at the international level gained momentum. In particular, the final quarter of the twentieth century saw growing recognition of the need for internationally coordinated protection strategies aimed at maintaining the Earth’s capacity to sustain development and support life. This period marked a shift from largely national conservation initiatives toward more globally aligned frameworks. (International Union for Conservation of Nature, 1980)

Despite the growing global uptake of protected areas, countries adopted diverse terminologies, governance models, and management approaches, resulting in substantial heterogeneity across protected area systems. In response to this diversity, international efforts to standardise protected area definitions began as early as the 1930s and were later formalised through the work of the International Union for Conservation of Nature (IUCN) from the 1960s onward. A key milestone in this process was the introduction of a ten-category classification system in 1978, based on management objectives. Recognising limitations in this framework, including ambiguity and inconsistent application, the IUCN initiated a review process that culminated in the adoption of the six-category system (Categories Ia–VI) in 1994. This revised system emphasised primary management objectives rather than national terminology or ownership structures. (Phillips, 2004)

In recent decades, the IUCN protected area category system has evolved beyond its original role as a shared conceptual framework. It has increasingly been embedded in national legislation, protected area system planning, and international policy processes. The endorsement of the category system by the Convention on Biological Diversity in 2004 represented a significant turning point, establishing it as a global reference standard for reporting and planning. At the same time, debates regarding activities such as resource extraction within protected areas have highlighted challenges in the interpretation and application of the system. These discussions led to the “Speaking a Common Language” project, which evaluated the effectiveness of the category framework and proposed revisions to improve consistency, transparency, and stakeholder engagement.

\begin{figure}[H]
    \centering
    \includegraphics[width=6.29966in,height=3.19444in]{Bildschirmfoto 2025-11-05 um 13.18.23.png}
    \label{fig:placeholder}
\end{figure}
\emph{Figure 1. Change in protected area and OECM coverage since 2000}

Recent data illustrate the rapid global expansion of protected areas over the past two decades. According to the United Nations Environment Programme World Conservation Monitoring Centre, a total of 302,934 protected areas have been reported to the Protected Planet Initiative, covering approximately 54.29 million km\textsuperscript{2} of the Earth’s surface. As of 2024, this corresponds to 17.6\% of terrestrial and inland waters and 8.4\% of marine and coastal areas being designated as protected globally. (UNEP-WCMC and IUCN, 2024)

\subsection{Biases in the Siting of Protected Areas} \label{biases-pa}

Assessing conservation progress solely on the basis of the total area covered by protected areas can be misleading. As Williams et al. (2022) argue, biodiversity conservation outcomes depend critically on where protected areas are established, particularly whether protection is concentrated in locations with high ecological value, such as regions hosting threatened species. Emphasising aggregate coverage targets without sufficient attention to spatial allocation risks undermining conservation objectives by neglecting factors such as ecosystem representation, connectivity, and effective management (Barnes, 2015).

Empirical evidence suggests that protected areas are unevenly distributed across the globe. Large protected area complexes are concentrated in regions such as the Amazon Basin, the Congo Basin, and parts of North America, Australia, and Russia, whereas areas including Southeast Asia, the Middle East, and North Africa exhibit comparatively limited coverage. These spatial disparities extend beyond geographic regions to differences across biomes. Tropical moist forests, particularly in South America and Central Africa, are relatively well protected, while tropical dry forests, grasslands, savannas, and wetlands remain underrepresented. Marine ecosystems are especially underprotected, with less than 1\% under formal protection. (Chape et al., 2008)

\begin{figure}[H]
    \centering
    \includegraphics[width=0.5\linewidth]{figures/Bildschirmfoto 2025-11-05 um 13.24.26.png}
    \label{fig:placeholder}
\end{figure}
\emph{Figure 2. Protected areas of the world, August 2025. Terrestrial protected areas are shown in green, marine protected areas in blue. Source: WDPA}

Beyond geographic and biome-level disparities, biases in protected area coverage also arise with respect to the protection of ecosystems that deliver critical services for human well-being. Areas of high ecological importance frequently provide essential ecosystem services such as clean freshwater provision and carbon sequestration (Watson et al., 2023). Historically, limitations in data availability and analytical methods have constrained the ability to systematically map these benefits, leaving global monitoring frameworks without agreed indicators to assess whether areas providing key ecosystem services are adequately protected (UNEP-WCMC and IUCN, 2024).

Recent advances have begun to address this gap. Chaplin-Kramer et al. (2022) introduce the concept of critical natural assets, defined as natural and semi-natural ecosystems that collectively deliver 90\% of the current magnitude of 14 types of nature’s contributions to people, including clean water, flood protection, fisheries, and carbon storage. In addition to their importance for human well-being, these critical assets also provide habitat for a substantial share of global biodiversity, covering 73\% of bird species and 66\% of mammal species. Despite their significance, only 15\% of these assets are formally protected, indicating substantial conservation gaps. Mapping critical natural assets may therefore play an important role in guiding future protected area establishment, also in the context of the 30×30 conservation target (Dudley and Stolton, 2022). These developments reflect an evolving understanding of protected areas, as growing demand for ecosystem services has prompted calls for new protected area categories aimed at strengthening ecological security for human populations (Xu et al., 2017).

An especially important ecosystem service provided by protected areas is climate change mitigation, a function that is expected to gain increasing prominence in global climate policy (Schirpke et al., 2021). Protected areas contribute to carbon storage through two distinct mechanisms. First, a development inhibition effect arises when the designation of protected areas directly constrains urban expansion and preserves existing carbon sinks. Second, a development replacement effect occurs when development pressures are displaced from ecologically valuable areas toward locations with lower environmental value (Hwang et al., 2021). However, the carbon sequestration potential of protected areas is itself vulnerable to climate change, as rising temperatures increase plant and soil respiration, partially offsetting the benefits of CO\textsubscript{2} fertilisation (Melillo et al., 2016). Maintaining and expanding well-managed protected areas is therefore critical not only for biodiversity conservation but also for climate change mitigation, reflecting the close interdependence between these objectives.

Empirical evidence from Brazil’s Amazon biome illustrates the climate mitigation role of protected areas. Nogueira et al. (2018) quantify carbon storage and deforestation-related carbon losses within protected areas in Amazonia, which contains approximately 20\% of the carbon stored in the world’s terrestrial vegetation. Their findings show that the majority of carbon stocks are located within Indigenous lands and federal and state conservation units, underscoring the importance of maintaining protected area integrity as a climate mitigation strategy. As the contribution of protected areas to carbon sequestration becomes increasingly salient, understanding their spatial distribution and carbon storage capacity has gained prominence. Lago-González et al. (2024) address this issue by employing machine learning methods to predict the carbon absorption capacity of protected areas under intensifying climate change impacts.

Finally, biases in protected area siting also manifest in substantial variation in the size and fragmentation of protected areas across regions. In South America, protected areas tend to be relatively large, with average sizes per country often exceeding 1,000 km\textsuperscript{2} (Salvio and Gomes, 2018). In contrast, protected areas in Europe are typically small and highly fragmented, frequently covering less than 1 km\textsuperscript{2}, which may limit their ecological effectiveness relative to larger, contiguous protected areas (Gaston et al., 2008).

\subsection{Predicting Protected Area Establishment: Evidence from Empirical Studies} \label{pa-prediction-empirical}

Halting the ongoing and unprecedented loss of biodiversity requires not only the expansion of protected areas but also their systematic and strategic planning. Given limited conservation resources, protected area expansion must integrate ecological, economic, and social considerations to prioritise locations where conservation outcomes can be achieved most effectively (Eken et al., 2014). This necessity implies explicit trade-offs between ecological objectives, economic opportunity costs, and social impacts. Despite their substantial potential to mitigate the negative effects of global environmental change on ecosystem services, protected areas often fall short of this potential. One reason is the spatial mismatch between areas that are rich in biodiversity and those that deliver the greatest ecosystem services (Schirpke et al., 2021). In addition, the establishment of protected areas can impose costs on nearby communities by constraining land use and limiting rural development opportunities (Llopis et al., 2021). Such impacts can arise even where conservation strategies are tailored to local contexts, seek to balance sustainable resource use with habitat protection, and enjoy general support among local residents (Hausner et al., 2021). To address these challenges and achieve socially sustainable conservation outcomes in biodiversity-rich regions, several studies emphasise the importance of incorporating socio-cultural ecosystem service assessments that capture how different groups perceive, use, and relate to ecosystems (Ndayizeye et al., 2020).

Despite the importance of strategic planning, the criteria guiding the establishment of new protected areas are often opaque, making PA designation difficult to predict. Although systematic, science-based frameworks for protected area selection exist, they are frequently applied inconsistently in practice, creating a gap between conservation theory and implementation that undermines the long-term effectiveness of protected area networks (Branquart et al., 2008). Empirical research suggests that, rather than being determined solely by conservation priorities, protected area establishment has been shaped by factors such as accessibility, economic viability, and historical institutional pathways (Xu and Pan, 2019). The relative importance of these factors varies across contexts. For marine protected areas, for example, site selection is influenced by both habitat quality and species connectivity, with connectivity often playing the more decisive role (Berglund et al., 2012). At the global level, protected area expansion has tended to follow opportunistic patterns, favouring areas characterised by isolation and low population density rather than areas optimised for environmental representativeness or biodiversity protection (Baldi et al., 2017).

While protected areas are widely regarded as a central instrument for reducing biodiversity loss, they have also been criticised for being disproportionately established on land that does not interfere with economically important human activities, such as fertile agricultural areas, rather than in locations with the highest biodiversity densities (Venter et al., 2018). A growing body of empirical evidence indicates that the siting of new protected areas reflects a combination of ecological value and socio-economic conditions, with economic context often exerting a stronger influence on location decisions than ecological considerations alone.

In practice, protected areas have frequently been established in locations that are ecologically relevant but associated with relatively low economic opportunity costs (Dasgupta et al., 2024). Global analyses show that protected areas are systematically biased toward remote regions, less arable lands, high elevations, and areas with steep or infertile terrain that are less valuable for agriculture and other productive uses (Mouillot et al., 2024). Mouillot et al. (2024) further demonstrate that protected area designations tend to cluster in locations that are “remote or of low value to agriculture and fisheries,” and are more prevalent in higher-income countries. This pattern reflects what is commonly described as the “residual” placement of protected areas, whereby conservation occurs preferentially in areas that least conflict with intensive human land use.

Country-specific studies reinforce this general pattern. In Australia, both public and private protected areas have been shown to occur predominantly in areas of lower land capability and at greater distances from cities, rather than in prime agricultural or urban regions (Sharma et al., 2025). Similarly, evidence from the United States indicates that distance from urban centres is a strong predictor of land being designated as protected, with the probability of protection increasing as distance from cities grows (Steif, n.d.). Earlier work by Joppa and Pfaff (2009) likewise finds that protected area networks are biased toward locations where protection is least likely to prevent land conversion, such as areas at higher elevations, with steeper slopes, and farther from roads and cities. Taken together, these findings suggest that economic factors, including land value and development pressure, play a central role in shaping where new protected areas can be successfully established.

Although ecological importance is a key consideration in conservation planning, empirical evidence indicates that high biodiversity alone is often insufficient to secure protected area designation when socio-economic conditions are unfavourable. While conservation planners tend to prioritise biodiversity-rich habitats, many areas with high conservation value remain unprotected due to economic pressures, governance constraints, or competing land-use demands (Sharma et al., 2025). 

Strict protected areas (IUCN Category I) are disproportionately concentrated in countries or regions characterised by high Human Development Index (HDI) scores and strong non-governmental organisation (NGO) presence. In such contexts, institutional capacity and political support facilitate the protection of ecologically valuable sites. In contrast, regions with high biodiversity but intense human pressure or strong dependence on natural resources tend to experience fewer new protected area designations. This trade-off between ecological value and socio-economic feasibility is  illustrated by spatial analyses that distinguish between potential and unrealistic conservation gains. Even within established biodiversity hotspots, protected area siting reflects these constraints. (Mouillot et al., 2024)  

Taken together, this body of evidence suggests that the spatial predictors of new protected area establishment span both environmental and economic dimensions. High conservation value, such as intact habitats, high vegetation indices, or elevated species richness, may be necessary to justify protection, but low economic opportunity costs, such as marginal land quality, sparse populations, or lower local economic activity, often make protection politically and economically feasible. Understanding this balance is critical for the success of the 30×30 target. If future protected area expansion continues to follow historical patterns that favour low-conflict locations, its capacity to halt biodiversity loss may remain limited. (Sharma et al., 2025)

A recent study by Greenhill et al. (2024) demonstrates how specific environmental and economic variables can be operationalised within predictive machine learning models of conservation decision-making. The authors analyse jurisdictional determinations related to the designation of protected water resources in the United States. Using aerial imagery and geophysical data, the authors train a deep learning model to predict where new protected areas are most likely to be established. While this study focuses primarily on water-related protected areas and emphasises geographical and ecological inputs, leaving economic factors largely unmodelled, it offers a valuable proof of concept for applying machine learning to predict protected area establishment. As such, it serves as an important methodological reference point for subsequent research.

Building on these insights, this thesis addresses key gaps in the existing literature by explicitly integrating economic and financial indicators into spatial machine learning models that predict new protected area designations under the 30×30 biodiversity target. Whereas prior approaches focus predominantly on ecological or socio-environmental variables and often rely on linear modelling frameworks (Buhler and Benson, 2024), this research incorporates subnational economic activity and transition-risk-relevant indicators alongside traditional ecological predictors. By combining these dimensions within non-linear machine learning models, the thesis aims to provide a more nuanced and policy-relevant forecasting framework for protected area expansion, reflecting both conservation priorities and the economic constraints that shape real-world decision-making.

\subsection{Machine-Learning Applications in Conservation} \label{ml-conservation}

Machine-learning methods are increasingly employed in conservation science to generate spatially explicit predictions. Existing applications are particularly prevalent in areas such as habitat suitability mapping, land-cover and land-use change detection, and ecological risk modelling. These approaches have demonstrated strong performance in integrating large volumes of remote-sensing, environmental, and geospatial data. However, the majority of current applications focus on predicting \textit{biophysical} outcomes, while comparatively few studies apply machine learning to forecast \textit{policy-driven} or institutional conservation decisions.

A growing body of empirical work illustrates the breadth of machine-learning applications in conservation contexts. Da Silveira (2021), for example, applies species-distribution modelling techniques to map coral habitat suitability within Brazilian marine protected areas, demonstrating the use of machine learning for fine-scale ecological prediction in marine systems. Doyle (2021) employs a Random Forest classifier trained on harmonised Landsat imagery to detect forest and wetland loss in Belize across multiple decades, highlighting how spectral, topographic, and environmental predictors can be combined to produce high-resolution spatial assessments of land-cover change.

Other studies extend machine-learning methods to model human-driven landscape dynamics. Mayfield (2017) compares statistical and machine-learning approaches to predict deforestation risk in Mexico and Madagascar, identifying accessibility, land-use dynamics, population pressure, and proximity to protected areas as key predictors. More recently, machine learning has also been applied to institutional classification tasks that go beyond purely ecological outcomes. Greenhill (2024), for instance, trains a ResNet-based deep learning model to predict whether waterways in the United States fall under federal jurisdiction, illustrating that geospatial machine-learning techniques can be successfully applied to regulatory and legal decision contexts.

Taken together, this literature demonstrates the flexibility and predictive power of machine-learning methods in conservation-related applications. At the same time, it reveals a strong orientation toward retrospective or contemporaneous classification tasks, with relatively limited attention to forecasting future conservation policy decisions.

\subsection{Synthesis and Research Gap} \label{synthesis-gap}

Several consistent patterns emerge from the existing literature on machine learning in conservation. First, predictive models rely heavily on remote-sensing data, terrain variables, and other biophysical inputs derived from Earth-observation systems. Second, a wide range of modelling approaches is employed, including both traditional statistical methods (such as logistic regression) and machine-learning techniques (such as Random Forests and deep neural networks), with model choice often driven by data availability and task complexity. Third, the majority of predictive applications focus on ecological or land-use outcomes, such as habitat suitability, deforestation risk, or ecosystem condition, rather than on the processes that govern conservation decision-making itself.

As a result, a clear research gap remains. Despite extensive evidence that protected area establishment is shaped jointly by ecological value, socio-economic conditions, and governance constraints, existing machine-learning studies rarely attempt to forecast \textit{future protected area designations}. Where conservation-related decisions are modelled, predictors are often limited to hand-crafted environmental variables, and models are typically not designed to capture the complex spatial, economic, and institutional factors relevant for policy-driven land-use change.

This thesis addresses this gap by developing a geospatial machine-learning framework that explicitly targets the prediction of protected area establishment. The approach integrates high-dimensional Earth-observation data with economic and governance indicators to estimate the probability that currently unprotected areas will be designated as protected in the future. By focusing on predictive rather than purely descriptive modelling, the framework aims to support more informed conservation planning and transition risk assessment.

More broadly, recent work highlights the potential for machine learning and artificial intelligence to enhance the identification, design, and adaptive management of protected areas by improving data integration, analysis, and decision-making processes (Mühling, 2023). This thesis contributes to this emerging agenda by extending machine-learning applications in conservation toward the prediction of policy-relevant conservation outcomes.
\newpage

\section{Data and Methodology}\label{data-and-methodology}

\subsection{Study Area and Spatial Framework} \label{Study-Area-and-Spatial-Framework}

\subsubsection{Geographic scope and justification} \label{Geographic scope and justification}

The study area includes the terrestrial territory of South America, including all continental land areas and major islands, while excluding open ocean. South America is chosen as the geographic focus for several reasons. First, the continent contains some of the world’s most biodiverse ecosystems, including the Amazon rainforest, the Andes, the Atlantic Forest, and extensive savanna and wetland systems, making it a central region for global biodiversity conservation efforts. Second, South America has experienced substantial expansion of protected areas over recent decades, providing sufficient temporal variation to study patterns of protected area establishment. Third, the region shows pronounced heterogeneity in ecological conditions, economic activity, infrastructure, and governance capacity, allowing the analysis to capture a wide range of drivers relevant for conservation decision-making.

Focusing on a single continent ensures institutional and ecological coherence while maintaining sufficient scale for meaningful spatial generalisation. This design choice avoids conflating fundamentally different legal and governance regimes across continents, while still providing a large and diverse empirical setting in which patterns of protected area expansion can be analysed systematically.

\subsubsection{Spatial resolution and grid construction} \label{Spatial resolution and grid construction}

All analyses are conducted on a regular spatial grid with a resolution of 1~km $\times$ 1~km. This grid serves as the common reference framework for integrating heterogeneous geospatial datasets, including environmental variables, land-use indicators, economic proxies, and protected area boundaries. The grid-based approach enables consistent spatial alignment across data sources with different native resolutions and supports pixel-level modelling of protected area establishment.

Each grid cell represents a spatial unit that is observed repeatedly over time, forming the basis for a pixel--year panel dataset used in subsequent analyses. A backbone layer with the exact terrestrial extent of South America ensures that predictions of protected area establishment are restricted to locations where protected areas could plausibly be designated.

The choice of a 1~km $\times$ 1~km spatial resolution reflects a trade-off between spatial detail, and the native resolution of the majority of input datasets. While higher-resolution land-cover products exist, many key predictors relevant for protected-area designation such as population density, climate variables, and long-run infrastructure layers are natively available only at coarser spatial scales. A 1~km grid therefore provides a consistent and information-preserving spatial support that enables integration across heterogeneous data sources while maintaining sufficient spatial granularity for continental-scale analysis.

\subsubsection{Coordinate reference system} \label{Coordinate reference system}

The reference grid is defined in the coordinate reference system EPSG:3857 (WGS 84 / Web Mercator). The reference grid is defined suitable for continental-scale spatial analysis in South America. All spatial datasets are reprojected to a common coordinate reference system prior to rasterisation or aggregation to ensure spatial consistency across inputs. The use of a single, consistent coordinate reference system prevents alignment errors when combining environmental, economic, and institutional data layers. While EPSG:3857 introduces known distortions in area at high latitudes, these effects are limited within the latitudinal extent of South America and are consistent across all inputs.

The use of EPSG:3857 ensures seamless integration with  derived datasets and avoids repeated reprojection and interpolation across data sources. Projection-induced distortions in area and distance are constant across the South American domain and therefore do not affect relative distance measures, spatial rankings, or the discriminative performance of the machine-learning models, which rely on comparative rather than absolute spatial relationships. The unit of observation remains 1km, which is the most relevant scale.

\subsection{Data}
\label{data}

This thesis integrates a diverse set of geospatial datasets to construct both the target variable (protected area establishment) and a comprehensive set of predictor variables capturing ecological value, land use, environmental conditions, economic activity, infrastructure, and governance-related proxies. The combination of these data sources is designed to reflect the multidimensional nature of protected area designation.

All datasets are harmonised to a common spatial framework and temporal structure prior to analysis, enabling pixel-level modelling of protected area establishment across South America. Unless stated otherwise, datasets are processed to match the 1~km $\times$ 1~km reference grid described in Section~\ref{Study-Area-and-Spatial-Framework}.

\subsubsection{Overview of Datasets}
\label{Overview-of-Datasets}

Table~1 provides an overview of the primary datasets used in this study, summarising their thematic focus, relevance for predicting protected area establishment, and data sources. The datasets can be grouped into three broad categories: (i) existing conservation context, (ii) environmental value and biophysical conditions, and (iii) economic activity, infrastructure, and governance-related proxies. This categorisation reflects the conceptual framework developed in the literature review, which emphasises that protected area siting is driven by the interaction of ecological importance and socio-economic feasibility.
\newline

\begin{longtable}{%
  >{\raggedright\arraybackslash}p{0.25\linewidth}%
  >{\raggedright\arraybackslash}p{0.45\linewidth}%
  >{\raggedright\arraybackslash}p{0.30\linewidth}%
}
\caption{Overview of datasets and variables used to predict protected area establishment.} \\
\toprule
\textbf{Data Category} & \textbf{Description and Role in the Analysis} & \textbf{Data Source and Coverage for Analysis} \\
\midrule
\endfirsthead

\toprule
\textbf{Data Category} & \textbf{Description and Role in the Analysis} & \textbf{Data Source and Coverage for Analysis} \\
\midrule
\endhead

\bottomrule
\endlastfoot

\multicolumn{3}{l}{\textbf{Existing Conservation Context}} \\
\addlinespace
Protected Areas (WDPA) &
Global database of terrestrial protected areas used to construct the binary protection indicator and to identify transitions from unprotected to protected status over time. &
World Database on Protected Areas (WDPA); terrestrial protected areas with designation year metadata \\

\addlinespace\addlinespace

\multicolumn{3}{l}{\textbf{Environmental and Ecological Conditions}} \\
\addlinespace
Vegetation index (NDVI) &
Normalized Difference Vegetation Index as a proxy for vegetation density and productivity. Used to characterise ecological conditions and land cover dynamics. &
MODIS NDVI, annual composites, 2000--2024, Google Earth Engine \\

Deforestation and forest cover &
Tree canopy cover in the year 2000 and annual forest loss thereafter. Captures forest extent, degradation, and land-use pressure over time. &
Hansen Global Forest Change (Landsat-based), 2000--2024, Google Earth Engine \\

Wildfire occurrence &
Annual burned-area indicator capturing exposure to fire-related disturbance and ecosystem stress. &
Global wildfire / burned-area products, 2000--2024, Google Earth Engine \\

Elevation and slope &
Topographic characteristics influencing land accessibility, agricultural suitability, and conservation costs. Slope computed from elevation prior to resampling. &
Shuttle Radar Topography Mission (SRTM), static (2008), Google Earth Engine \\

Climate variables &
Bioclimatic variables capturing long-run temperature and precipitation patterns, seasonality, and extremes. Used to characterise environmental suitability and constraints. &
WorldClim bioclimatic variables, static climatology, Google Earth Engine \\

Global Safety Net (GSN) layers &
Spatial indicators of globally important conservation areas, including: \begin{itemize}
    \item High biodiversity areas
    \item Intact wilderness areas
    \item Climate stabilisation areas
    \item Potential wildlife corridors
    \item Terrestrial ecoregions
\end{itemize}
Used as high-level ecological context variables. &
Global Safety Net dataset (peer-reviewed), static, 2020 \\

\addlinespace\addlinespace

\multicolumn{3}{l}{\textbf{Human Pressure, Economic Activity, and Infrastructure}} \\
\addlinespace
Population density &
Gridded population counts capturing human presence and land-use pressure. Used as a proxy for development intensity and opportunity costs of protection. &
Gridded Population of the World (GPW), 5-year intervals, 2000--2020, Google Earth Engine \\

Night-time lights (HNTL) &
Harmonized night-time light intensity capturing economic activity and built infrastructure while correcting for sensor discontinuities. &
Harmonized Night-Time Lights (HNTL), annual, 2000--2024 \\

Road infrastructure &
Major road network, capturing accessibility and development pressure. &
Global Roads Inventory Project (GRIP), static, 2018 \\

Energy and extractive infrastructure &
Power plants and oil and gas infrastructure as proxies for industrial activity and fixed capital investment. &
Global Power Plant Database (WRI, 2018); Oil and Gas Infrastructure Mapping (EDF, 2023), Google Earth Engine \\

\end{longtable}

\subsubsection{Target Variable: New Protected-Area Designations}
\label{target-variable}

The target variable indicates whether a grid cell transitions from an unprotected to a protected status within a fixed forward-looking time window. Protected-area boundaries and designation dates are obtained from the World Database on Protected Areas (WDPA), which provides spatial polygons with associated temporal metadata.

WDPA polygons are rasterised to the 1~km $\times$ 1~km reference grid for each year of observation. A grid cell is classified as protected in a given year if any portion of the cell overlaps with a protected area designated as active in that year. Cells that are already protected at the start of the observation period are excluded from the risk set, such that the analysis focuses exclusively on locations eligible for new protected-area designation.

For each pixel--year observation, the target variable is defined as a binary indicator equal to one if the cell becomes protected at any point within the subsequent \textit{N}-year lookahead window, and zero otherwise. Observations for which the full lookahead window cannot be observed are right-censored and excluded from model training and evaluation to ensure consistent label definition.

The target captures the initial transition into protected status. Subsequent disestablishment, boundary contraction, or reclassification events are not modelled explicitly; pixels are treated as having experienced a successful transition at the time of first designation. The analysis therefore focuses on the determinants of initial protected-area establishment rather than long-run legal persistence.

\subsubsection{Predictor Variables}
\label{predictor-variables}

Predictor variables are designed to capture both the ecological suitability of locations for conservation and the socio-economic and institutional factors that shape protected area establishment decisions. Consistent with the conceptual framework developed in the literature review, predictors are grouped into three broad categories reflecting existing conservation context, environmental and ecological conditions, and human pressure and economic activity.

\paragraph{Existing conservation context.}

This group of predictors characterises the surrounding conservation landscape and captures spatial dependence in protected area expansion. Variables include indicators of existing protected areas and the distance to the nearest previously established protected area. These measures reflect path dependence and clustering effects, whereby new protected areas are more likely to be designated near existing ones due to administrative, ecological, or political considerations.

\paragraph{Environmental value and ecosystem characteristics.}

Environmental predictors describe the biophysical attributes of each grid cell that are relevant for biodiversity conservation. These include vegetation indices, forest cover and forest loss, wildfire occurrence, climate variables, topographic characteristics such as elevation and slope, and global conservation priority layers. Together, these variables capture ecosystem condition, habitat suitability, environmental stressors, and long-run climatic constraints that influence the conservation value of a location.

\paragraph{Economic activity, infrastructure, and governance.}

This set of predictors represents human pressure and opportunity costs associated with land protection. Variables include population density, night-time light intensity as a proxy for economic activity, and distances to major infrastructure such as roads, power plants, and oil and gas facilities. These indicators reflect accessibility, development intensity, and fixed capital investments that may raise the economic or political cost of establishing protected areas.

\paragraph{Static versus dynamic predictors.}

Predictor variables are classified as either static or dynamic depending on their temporal variation. Static predictors, such as elevation, slope, and long-run climate variables are assumed time-invariant and repeated across all years for a given grid cell. Road- and energy infrastructure layers are also treated as static due to limited historical coverage; this may attenuate short-run effects of newly constructed infrastructure. Dynamic predictors, including vegetation indices, deforestation, wildfire occurrence, population density, and night-time lights, vary over time and are aligned to the year of observation.

\paragraph{Treatment of missing values.}

Missing values arise primarily from genuine data limitations, such as incomplete temporal coverage in specific remote-sensing products. No global temporal imputation is applied. Instead, missing values are retained explicitly as \texttt{NA} and handled natively by the tree-based machine-learning models employed in this study, which can incorporate missingness through their internal split logic.

For selected dynamic predictors with coarse or partially incomplete temporal coverage, limited forward-filling is applied within each grid cell under the assumption of gradual temporal change. Specifically, gridded population density (GPW), available at five-year intervals, is forward-filled between observation years to ensure annual coverage, and harmonized night-time lights (HNTL) are forward-filled over a short gap at the end of the series where annual observations are unavailable. All remaining missing values are left unfilled and treated natively by the learning algorithms. This strategy avoids introducing artificial variation through aggressive imputation while preserving maximal data coverage.

\subsection{Spatial Harmonisation and Panel Construction}
\label{Spatial-Harmonisation}

All datasets are harmonised to a common spatial and temporal framework prior to model estimation. Given the heterogeneous origins of the input data , this harmonisation step is essential to ensure consistency and reproducibility. The harmonisation process consists of reprojection and resampling, spatial masking to a land-only backbone, temporal aggregation, and the construction of a unified pixel--year panel dataset.

\subsubsection{Reprojection and resampling}
\label{resampling}

All spatial datasets are harmonised to a common projected reference grid defined in EPSG:3857 with a spatial resolution of 1~km $\times$ 1~km. This grid is instantiated as a fixed land-only backbone covering South America and serves as the single authoritative spatial reference throughout the pipeline.

Where possible, reprojection and resampling are performed at the data-extraction stage (e.g., in Google Earth Engine). All raster exports are generated directly in EPSG:3857 at a scale of 1,000~m. Continuous variables, such as climate, terrain, population density, and night-time lights, are resampled using bilinear interpolation to reduce artefacts from scale transformation. Categorical or binary datasets, including land-cover classes, protected-area indicators, and rasterised infrastructure layers, are resampled using nearest-neighbour interpolation to preserve discrete class boundaries.

To ensure consistency across all layers and time periods, a fixed land-only backbone grid is constructed. This backbone defines the universe of grid cells eligible for analysis and is used to mask all datasets uniformly. Only pixels corresponding to terrestrial land areas within South America are retained, while ocean pixels and areas outside the study region are excluded.

During preprocessing and panel construction in Python, all raster layers are systematically checked for alignment with the backbone grid in terms of coordinate reference system, resolution, and spatial extent. Any raster that does not match the backbone grid exactly is reprojected and resampled to the reference grid, using nearest-neighbour resampling to avoid introducing additional spatial smoothing.

\subsubsection{Temporal aggregation}
\label{temporal-aggregation}

Input datasets differ in temporal resolution and sampling frequency. To ensure consistency with the annual structure of protected-area designation data, all time-varying predictors are represented at the calendar-year level. Temporal handling occurs in two stages: within-year aggregation during data extraction in Google Earth Engine (GEE) and year-level assembly of a complete pixel--year panel during merging.

For image collections with multiple observations per year, a single annual raster is constructed per calendar year. Continuous remote-sensing indicators (e.g., NDVI, VIIRS night-time lights, Dynamic World class probabilities) are aggregated using annual medians to reduce sensitivity to outliers and transient artefacts. Wildfire exposure is aggregated cumulatively using annual sums and derived binary indicators capturing whether burning occurred during the year. Variables provided as year-specific products (e.g., MODIS land cover) are used directly without within-year reduction. Annual and cumulative deforestation measures are derived from the Hansen loss-year layer. Protected-area status is constructed cumulatively by year by rasterising all WDPA polygons designated up to the observation year. Population density and satellite embeddings are exported as single images for each available year. Static layers, such as topography and long-run climate variables, have no temporal aggregation.

\subsubsection{Panel dataset construction}
\label{panel-construction}

Following spatial alignment and temporal harmonisation of all predictors, the datasets are merged into a unified pixel--year panel. Each observation corresponds to a single grid cell in a specific calendar year and contains three components: (i) the binary target variable indicating future protected area designation within the specified lookahead window; (ii) the full set of static and dynamic predictor variables; and (iii) unique identifiers for the spatial cell and year.

The final panel consists of approximately $20.8$ million unique land pixels observed over $25$ years (2000--2024). Since pixels already protected at the start of the panel are excluded, as stated in section~\ref{target-variable}, there are a total of $413{,}722{,}039$ pixel--year observations. Transitions into protected status are rare events: fewer than $0.5\%$ of eligible pixels transition in any given year. This extreme class imbalance is a defining characteristic of the data and directly motivates the modelling choices discussed in Section~\ref{framework}.

\subsection{Feature Engineering and Multi-Scale Representation}
\label{features}

The pixel--year panel integrates a comprehensive set of biophysical, socio-economic, and anthropogenic predictors. To reflect the fact that protected-area designation is shaped not only by conditions at a specific location but also by its surrounding landscape, the feature set combines (i) baseline cell-level characteristics with (ii) explicitly engineered multi-scale spatial context features.

This design allows the modelling framework to capture both local suitability for conservation and broader spatial patterns related to ecological connectivity, development pressure, and coordinated land-use planning.

\subsubsection{Baseline Cell-Level Features}
\label{baseline_features}

For every grid cell and year, the panel includes the full set of harmonised predictor variables described in Section~\ref{predictor-variables}. These baseline features characterise conditions at the pixel itself and represent the immediate local environment from which transitions into protection may arise.

In event-history terms, the baseline features describe the contemporaneous ``state'' of each pixel in a given year, conditional on not yet being protected. They form the core explanatory variables linking local conditions to the probability of future protected-area designation.

\subsubsection{Multi-Scale Spatial Smoothing}
\label{scale}

Protected-area establishment decisions are influenced by spatial processes operating at multiple scales. While some drivers act locally, others reflect broader regional dynamics such as agricultural frontiers, infrastructure corridors, or clusters of high ecological value.

To incorporate these neighbourhood effects directly into the dataset, spatially smoothed representations are computed for a selected subset of dynamic predictors. For each pixel-year observation, uniform averaging windows of size $4\times4$, $16\times16$, and $64\times64$ grid cells are applied to variables including NDVI, population density, night-time lights, deforestation, and wildfire activity. These window sizes correspond to neighbourhood extents of 4\,km, 16\,km, and 64\,km, respectively. The window sizes were chosen to span local, regional, and sub-continental scales; alternative window definitions were explored in preliminary analyses and yielded qualitatively similar results.

The selected window sizes (4~km, 16~km, and 64~km) are designed to capture spatial processes operating at local, regional, and landscape scales, respectively. Smaller windows reflect immediate neighborhood effects such as local disturbance and accessibility, while larger windows capture broader land-use patterns, ecological gradients, and coordinated planning processes that extend beyond individual grid cells. This multi-scale representation allows the model to incorporate spatial context without imposing hard administrative boundaries or requiring computationally intensive hierarchical architectures.

\subsubsection{Distance-Based Features}
\label{distance}

In addition to neighbourhood-level smoothing, the panel includes distance-based features that quantify each pixel’s proximity to existing protected areas and infrastructure. For each year, the Euclidean distance to the nearest protected-area boundary (\texttt{dist\_wdpa}) is computed on the reference grid. Analogous distance measures are constructed for road networks and energy infrastructure.

These features capture spatial dependence in conservation outcomes, reflecting processes such as spatial clustering of protected areas, buffer-zone formation, and incremental expansion adjacent to existing reserves. By encoding proximity explicitly, distance-based measures complement local averages and allow the model to distinguish between contextual neighbourhood effects and direct spatial adjacency.

\subsection{Modelling Framework}
\label{framework}

This section describes the modelling strategy used to predict the emergence of new protected areas. The approach combines a temporally explicit transition formulation with validation procedures designed to ensure temporal honesty and spatial robustness.

\subsubsection{Problem Formulation and Target Definition}
\label{problem-formulation}

The objective is to predict whether a currently unprotected grid cell will be designated as protected within a fixed future horizon. Let $y_{i,t}$ indicate whether pixel $i$ is protected in year $t$. The target variable is defined as a five-year forward-looking transition indicator:

\[
\text{transition}_{01\_win5,i,t}
=
\begin{cases}
1 & \text{if } y_{i,t}=0 \ \text{and} \ y_{i,t+k}=1 \ \text{for any } k \in \{1,\dots,5\}, \\
0 & \text{otherwise}.
\end{cases}
\]

This formulation captures whether protection occurs within a policy-relevant
five-year window, consistent with medium-term conservation planning horizons. All models are trained exclusively on the \emph{risk set} of unprotected pixels ($y_{i,t}=0$), ensuring that predicted probabilities correspond to well-defined transition risks. The resulting classification problem is highly imbalanced, with positive events representing approximately $0.3$--$0.5\%$ of eligible pixel--year observations.

\subsubsection{Relation to Discrete-Time Hazard Models}
\label{discrete-time}

Conceptually, the modelling framework is closely related to discrete-time hazard (or duration) models commonly used in econometrics to analyse the timing of events. In this interpretation, each unprotected pixel is ``at risk'' of transition, and the model estimates the hazard of PA designation as a function of time-varying and time-invariant covariates.

Unlike traditional parametric hazard models, which require strong assumptions about baseline hazard functions and covariate effects, the machine-learning approaches employed here allow for flexible, nonlinear relationships and complex interactions between predictors. The use of tree-based ensemble models therefore provides a data-driven approximation of the hazard function without imposing functional-form restrictions. Importantly, the discrete-time formulation preserves temporal ordering and allows covariates to evolve over time, aligning the modelling strategy with the panel structure of the data and the institutional reality of annual PA designation
decisions.

\subsubsection{Rare-Event Prediction Challenges}
\label{challenges}

Predicting protected-area establishment poses several challenges typical of rare-event settings. First, the extreme class imbalance can bias standard classification algorithms toward predicting non-transition outcomes, reducing their usefulness for identifying future PA expansions. Second, small absolute changes in predicted probabilities can correspond to large differences in ranking performance, which is critical for spatial prioritisation tasks. 

To address these issues, the modelling framework focuses on probabilistic outputs and ranking-based evaluation metrics rather than accuracy-based measures. Class imbalance is handled through a combination of class weighting and algorithm-specific parameters designed to stabilise learning under skewed outcome distributions. No synthetic resampling is applied, ensuring that the empirical distribution of transitions is preserved.

\subsection{Machine-Learning Models}
\label{models}

This study employs a set of machine-learning models to predict the establishment of new protected areas (PAs) under the global 30$\times$30 biodiversity target. The modelling framework is formulated as a forward-looking transition-risk problem and is evaluated across two geographic scopes: (i) a continental-scale model for South America (Model~1), which constitutes the main empirical analysis, and (ii) a country-level model for Colombia (Model~C), used for methodological development and robustness checks. Across these variants, three modelling approaches are considered: a static suitability baseline, a Random Forest transition model, and a LightGBM transition model. The LightGBM model serves as the primary production model for forecasting and spatial risk mapping.


\subsubsection{Static Suitability Baseline}
\label{static}

As a benchmark, a static suitability model is estimated using cross-sectional data from the year 2000. The model is trained to distinguish protected from unprotected pixels based solely on spatial characteristics observed at baseline. Suitability scores for pixels unprotected in 2000 are then compared to observed protected-area designations between 2000 and 2024. Evaluation of the static model uses the same right-censoring logic as the transition models to ensure comparability. This model captures long-run spatial correlates of protection, such as remoteness, topography, and baseline ecological value, but ignores temporal dynamics and policy evolution. It therefore establishes a lower bound on predictive performance and quantifies how much explanatory power arises from static spatial suitability alone.

\subsubsection{Random Forest Transition Model}
\label{RF}

The Random Forest transition model introduces temporal dynamics by directly predicting the five-year transition indicator. A Random Forest classifier is trained on the full risk set, with class imbalance addressed using per-tree class weighting (\texttt{class\_weight="balanced\_subsample"}). Random Forests provide a flexible, nonparametric benchmark capable of capturing nonlinearities and interactions among predictors. While more memory-intensive than gradient-boosting methods, this model offers useful diagnostic insights through feature importance measures and serves as an interpretable ensemble baseline.

\subsubsection{LightGBM Transition Model}
\label{LGBM}

LightGBM is used due to its computational efficiency and strong performance on large, imbalanced tabular datasets. The model is trained on the same transition target and temporal regime as the Randomm Forest model, using gradient-boosted decision trees.

Class imbalance is handled via the \texttt{scale\_pos\_weight} parameter, computed as the ratio of negative to positive observations in the training data. Optional time-based sample weights, increasing linearly from early to later years, are used in selected specifications to place greater emphasis on recent observations.

Hyperparameters governing tree structure and regularisation (including \texttt{num\_leaves}, \texttt{max\_depth}, and \texttt{min\_child\_samples}) are selected using temporal cross-validation and optimised for Precision--Recall performance. Final models are trained without early stopping, using a fixed number of boosting rounds derived from cross-validation to avoid information leakage.

\subsubsection{Separate Colombia Pipeline}
\label{colombia}

In addition to the continental-scale model (Model~1), a fully separate country-level pipeline is implemented for Colombia (Model~C). This pipeline mirrors the complete data processing, feature engineering, modelling, and evaluation workflow used at the continental scale, but is applied to a substantially smaller spatial domain. Both the Random Forest and LightGBM transition models are estimated using identical specifications and target definitions, with adaptations limited strictly to the geographic extent of the data.

The Colombia pipeline serves two complementary purposes. First, it provides a computationally tractable environment for rapid iteration, debugging, and methodological development under realistic data conditions. This includes experimentation with feature construction, probability calibration strategies, validation design, and diagnostic analyses that would be prohibitively expensive to perform repeatedly at the continental scale. Second, the pipeline functions as a robustness and transferability check, allowing assessment of whether modelling choices developed at smaller scale generalise coherently when applied to the full South American dataset.

Colombia is selected as the country-level case study because it exhibits substantial heterogeneity in ecological conditions, land-use patterns, infrastructure development, and conservation dynamics, encompassing many of the key drivers relevant for protected-area establishment across South America. While not intended to be fully representative of the continent, Colombia provides a rich and diverse empirical setting in which modelling assumptions and design choices can be stress-tested before large-scale deployment.

All final modelling decisions reported for the South America analysis are validated for internal consistency within the Colombia pipeline prior to full-scale estimation.
\subsubsection{Model Development and Tuning Procedure}
\label{model-development}

Note: Hyperparameter tuning via Optuna (15 trials, 5-fold rolling temporal CV) selected $n_{\text{estimators}} = 777$ for South
  America and n$_{\text{estimators}}$ = 847 for the USA. Both values were reduced to 300 for the final models. This decision was
   validated empirically: across all tested values (200, 300, 777/847), top-decile predictive accuracy (P@1%) was
  identical — 0.5386 for South America and 0.0400 for the USA regardless of n$_{\text{estimators}}$ — consistent with the
  well-established result that Random Forest performance saturates well before the Optuna-selected counts in
  large-sample settings (Oshiro et al., 2012; Probst \& Boulesteix, 2017). The reduction from ≈800 to 300 trees
  (roughly 60%) cut wall-clock training time by approximately the same fraction with no measurable loss in test
  performance. A value of 300 falls within the range commonly reported for continental-scale geospatial RF models.

  Suggested references for note 2:
  - Oshiro, T.M. et al. (2012). How many trees in a random forest? MLDM.
  - Probst, P. \& Boulesteix, A.L. (2017). To tune or not to tune the number of trees? JMLR.


Model development follows a structured, stepwise procedure designed to maximise predictive performance while maintaining transparency, reproducibility, and temporal integrity.

Hyperparameter tuning is conducted using temporal cross-validation and early- stopping sets that strictly precede the test period. For LightGBM, tuning focuses on tree complexity, learning rate, and regularisation parameters, with average precision as the primary optimisation criterion. The final number of boosting iterations is determined from cross-validation results and fixed prior to final training to avoid test-set leakage.

Out-of-fold predictions from temporal cross-validation are retained and used for probability calibration. Final models are then trained on the full training and validation period without early stopping and evaluated on the held-out test period.

\subsubsection{Probability Calibration}

To ensure that predicted probabilities are interpretable as transition risks, model outputs are calibrated using out-of-fold predictions from temporal
cross-validation. Calibration methods include Platt scaling and, in robustness checks, isotonic regression. The fitted calibrator is applied to held-out test data and future prediction periods, with probability clipping used to ensure numerical stability. All reported probability surfaces and risk maps are based on calibrated
LightGBM predictions.

Because class-weighting schemes used to address extreme imbalance modify the effective loss function, raw model outputs do not directly correspond to well-calibrated transition probabilities. Probability calibration is therefore a necessary step for meaningful risk interpretation. All reported probabilities are post-calibrated using out-of-fold predictions from temporal cross-validation, and only calibrated values are interpreted as estimates of transition risk and used for spatial risk mapping.

\subsection{Training, Validation, and Test Design}
\label{train-test}

This section describes the temporal structure used for model training, validation, and evaluation. Because protected-area designation is a forward-looking policy event, model assessment must preserve the temporal ordering of observations and avoid any leakage of future information into the training process.

\subsubsection{Temporal split strategy}
\label{split}

The panel is partitioned using strictly forward-looking temporal splits. All models are trained only on information available prior to the test period, reflecting the real-world forecasting task faced by policymakers and risk analysts.

The temporal structure is as follows:
\begin{itemize}
    \item \textbf{Training period:} 2000--2016
    \item \textbf{Early-stopping / validation period:} 2014--2016
    \item \textbf{Test period:} 2017--2019
\end{itemize}

The test period is capped at 2019 to ensure that all observations admit a full five-year lookahead window, avoiding right-censoring of the target variable. The last test year 2019 therefore includes PA establishments until 2024. Earlier years (2000--2016) are used exclusively for model fitting and hyperparameter selection. No observations from the test period influence model estimation, tuning, or calibration. This design ensures temporal honesty: performance metrics reflect the ability of the model to anticipate future protected-area designations based solely on past information.

\subsubsection{Robustness to alternative test periods}
\label{robsutness}
To assess the sensitivity of model performance to the choice of evaluation window, alternative test splits were explored in which the held-out period begins earlier (e.g. 2016–2019 and 2018–2019), while preserving the same five-year lookahead definition and right-censoring logic. Predictive performance and ranking quality were found to be qualitatively stable across these alternative splits. The final test period of 2017–2019 is adopted as the primary evaluation window because it provides a balanced trade-off between a sufficient sample size of observed transitions, while maintaining a large enough test set to confidently evaluate performance metrics.

\subsubsection{Lookahead window definition}
\label{lookahead}

As described in section~\ref{problem-formulation}, the prediction target is defined as a five-year forward-looking transition indicator. For each pixel--year observation $(i,t)$, the outcome records whether the pixel becomes protected at any point during the interval $[t+1, t+5]$. To guarantee that this target is well defined, the sample is right censored at \texttt{LAST\_LABEL\_YEAR = 2019}. Observations beyond this year are excluded from training and evaluation, even if predictor data are available. This restriction prevents mechanically mislabeling late observations as non-transitions simply because future information is not yet observed.

\subsubsection{Avoidance of temporal leakage}
\label{leakage}

Multiple safeguards are implemented to prevent temporal leakage. First, all predictors are constructed using only information available at or prior to the prediction year, and no feature incorporates forward-looking data. In particular, distance-based measures to protected areas are computed using a one-year lag, ensuring that the model cannot indirectly infer the location of newly designated protected areas within the same year. Variables that directly encode contemporaneous protected-area status, such as current protection indicators or boundary flags, are excluded from the feature set.

Second, early-stopping and hyperparameter tuning are performed using validation data that strictly precede the test period. Final production models are trained without early stopping, using a fixed number of boosting iterations determined from temporal cross-validation, ensuring that no information from the test set influences model complexity.

Third, probability calibration is based exclusively on out-of-fold predictions from temporal cross-validation or early-stopping data and is applied unchanged to the test period.

\subsubsection{Spatial generalisation check}
\label{spatial generalisation}

Although the primary evaluation dimension is temporal, spatial autocorrelation in environmental and socio-economic variables raises the risk that models may overfit region-specific patterns. To assess spatial robustness without violating the temporal structure of the panel, a spatial generalisation check is conducted within the held-out test period. Test observations are grouped into a small number of large, contiguous spatial regions (e.g.\ macro-regions or coarse spatial clusters), and performance metrics are computed separately for each group.

This diagnostic verifies that predictive performance is not driven solely by a subset of regions and that the model captures continent-scale drivers of protected-area expansion rather than narrowly localised spatial correlations. Unlike full spatial block cross-validation, this approach preserves the event-history structure of the data and remains computationally tractable at continental scale.

\subsubsection{Class imbalance handling}
\label{imbalance}

Transitions into protected status are rare events, representing approximately 0.5\% of eligible pixel--year observations. To address this imbalance without distorting the spatial or temporal structure of the data, no resampling or undersampling is applied in the main models. Instead, imbalance is handled through algorithm-specific weighting schemes. For LightGBM, the \texttt{scale\_pos\_weight} parameter is set equal to the ratio of negative to positive observations in the training data, computed separately for each temporal split. Random Forest models use per-tree class balancing via \texttt{class\_weight="balanced\_subsample"}.

This strategy preserves the true prevalence of protection events while allowing the models to focus learning capacity on the rare transition class. Model performance is evaluated using metrics appropriate for imbalanced settings, as described in Section~\ref{evaluation-metrics}.

\subsection{Evaluation Metrics and Model Diagnostics}
\label{evaluation-metrics}

Because transitions into protected status are extremely rare, conventional threshold-based metrics such as accuracy or the F$_1$-score are uninformative and potentially misleading. Model evaluation therefore focuses on metrics that (i) are robust to severe class imbalance and (ii) directly reflect the ranking quality of predicted risks, which is the primary object of interest for spatial prioritisation and transition-risk assessment. Model performance is assessed using the following complementary metrics:

\begin{itemize}
    \item \textbf{Precision--Recall Area Under the Curve (PR--AUC).}  
    PR--AUC summarises the trade-off between precision and recall across all classification thresholds and is the primary evaluation metric in this study. In rare-event settings, PR--AUC is sensitive to performance on the positive (transition) class and therefore provides a meaningful measure of discriminative power where accuracy- or ROC-based metrics can be misleading.

    \item \textbf{Receiver Operating Characteristic Area Under the Curve (ROC--AUC).}  
    ROC--AUC is reported as a complementary metric to facilitate comparison with existing studies and standard machine-learning benchmarks. While less informative under extreme class imbalance, ROC--AUC captures overall ranking performance across both classes and serves as a useful diagnostic for detecting degenerate or unstable models.

    \item \textbf{Precision at Top-$k\%$ and Lift (P@$k$, Lift@$k$).}  
    Precision is computed among the highest-ranked predictions (top 1\%, 5\%, and 10\% of pixels by predicted probability), reflecting the practical use case of identifying a limited set of locations most likely to be designated as protected in the near future. To account for differences in base rates across evaluation splits, precision is additionally reported as lift relative to random allocation. Lift explicitly normalises P@$k$ by the underlying positive rate, ensuring that high precision is not mechanically driven by changes in test-set size or class prevalence.

    \item \textbf{Region-specific performance.}  
    To assess spatial robustness, PR--AUC, P@$k$, and Lift@$k$ are computed separately within each macro-region used in the spatial generalisation check. This ensures that aggregate performance is not driven by a small subset of regions and that the model generalises across heterogeneous geographic contexts.

    \item \textbf{Error decomposition and diagnostic analysis.}  
    False positives and false negatives are analysed explicitly. False positives indicate areas with high predicted risk that are not designated within the evaluation window, which may reflect delayed policy action or unobserved institutional constraints. False negatives highlight missed transitions and potential blind spots in the predictor set.
\end{itemize}

In addition to discrimination metrics, \textbf{probability calibration} is assessed to ensure interpretability of predicted risks. Calibration curves and Brier scores are reported for the LightGBM model using calibrated probabilities derived from temporal cross-validation. This step is essential for interpreting predicted values as approximate transition probabilities rather than arbitrary scores.

Model interpretability is supported through feature-importance measures derived from both Random Forest and LightGBM models. These diagnostics are used to identify the dominant biophysical, socio-economic, and spatial drivers of protected-area establishment and to verify that model behaviour aligns with theoretical expectations.

\subsection{Data Quality Checks and Limitations}
\label{limitations}

Given the scale and heterogeneity of the data used in this study, systematic data quality checks are essential. At the same time, several limitations arise from data availability, measurement uncertainty, and the observational nature of the analysis. These limitations are explicitly acknowledged to ensure correct
interpretation of the results.

\paragraph{Spatial and temporal consistency.}
All raster and vector datasets are aligned to a common 1\,km reference grid and checked for consistency in spatial extent, resolution, and coordinate reference system prior to panel construction. Automated alignment checks verify that all inputs match the backbone grid in shape, transform, and CRS. Temporal coverage is validated to ensure that each predictor is available (or explicitly missing) for the corresponding year in the panel. Results may be sensitive to the chosen spatial resolution (modifiable areal unit problem). While 1 km represents a reasonable compromise between detail and data availability, finer or coarser grids could yield different spatial patterns.

\paragraph{Protected-area boundary and dating uncertainty.}
The World Database on Protected Areas (WDPA) represents the most comprehensive global source of protected-area information, but it is subject to known limitations. In particular, designation years may be missing, rounded, or inaccurately reported for some polygons, and updates may lag behind formal designation. Residual dating uncertainty may affect the precise timing of observed transitions. This limitation is mitigated by the use of multi-year lookahead windows, which reduce sensitivity to exact designation
dates.

\paragraph{Measurement error in predictors.}
Several predictors are proxies rather than direct measures of underlying
concepts. For example, night-time lights approximate economic activity, population density proxies human pressure, and NDVI captures vegetation greenness rather than biodiversity or other nature-related data per se. Measurement error in predictors can attenuate estimated relationships and reduce predictive performance in specific contexts.

\paragraph{Incomplete temporal coverage.}
Some datasets exhibit incomplete temporal coverage or sporadic gaps (e.g.\ night-time lights, fire products, or biodiversity layers). While selected variables are forward-filled as described in section~\ref{predictor-variables}, missing values arising from genuine data limitations are retained as \texttt{NA} rather than imputed. Tree-based models naturally accommodate missingness through split rules, but predictive uncertainty may increase in regions or periods with systematically lower data availability, even though such regions have not been observed in the study area.

\paragraph{Non-causal interpretation.}
The modelling framework is explicitly predictive rather than causal. Estimated relationships reflect patterns associated with historical protected-area designation decisions and should not be interpreted as causal effects of individual predictors. In particular, socio-economic variables may capture both direct influences on policy decisions and correlated institutional or political factors that are not directly observed.

\paragraph{Static treatment of infrastructure.} Road networks and energy infrastructure are treated as static layers due to the lack of consistent, high-resolution historical datasets covering the full study period. This implies that infrastructure present in later years may be implicitly visible to the model in earlier periods, introducing a potential look-ahead bias. This limitation is mitigated by the model’s emphasis on relative spatial ranking rather than causal interpretation.

\paragraph{External validity and future policy shifts.}
Model performance is evaluated under the assumption that future protected-area expansion follows patterns similar to those observed in the historical data. Structural changes in conservation policy, governance, or international financing mechanisms could alter these patterns. Consequently, forecasts should be interpreted as conditional on past institutional dynamics rather than as unconditional predictions.

\subsection{Reproducibility and Implementation Details}
\label{reproducibility}

The entire analysis is implemented using a modular, script-based workflow designed to ensure reproducibility and transparency.

\paragraph{Software and libraries.}
All data processing and modelling steps are implemented in Python. Key libraries include \texttt{numpy}, \texttt{pandas}, \texttt{pyarrow}, and \texttt{duckdb} for data handling; \texttt{rasterio}, \texttt{geopandas}, and \texttt{shapely} for geospatial processing; and \texttt{scikit-learn} and \texttt{lightgbm} for machine-learning models. Data extraction from Google Earth Engine is performed using the GEE Python API.

\paragraph{Computational environment.}
Large-scale preprocessing and model training are executed on a high-performance computing (HPC) cluster. Jobs are submitted via a workload manager and configured to use fixed random seeds where applicable to ensure reproducibility. Memory- efficient, streaming-based data loading is employed to handle the full continental panel without loading all observations into memory simultaneously.

\paragraph{Pipeline structure and data versioning.}
The workflow is organised into sequential stages: data extraction, spatial harmonisation, panel construction, feature engineering, split construction, hyperparameter training, model training, evaluation, and visualisations. Intermediate outputs are stored in columnar \texttt{Parquet} format with explicit schema definitions. Input datasets are versioned and checksum- verified to ensure consistency across runs. Model configurations and hyperparameters are stored in external configuration files to decouple code and experiment settings.

\paragraph{Randomness and determinism.}
Random components in model training (e.g.\ tree construction, subsampling) are controlled via fixed random seeds. For LightGBM, the final production model is trained with a fixed number of boosting iterations determined from temporal cross- validation, avoiding stochastic early stopping in the final run.

\paragraph{Outputs and replication.}
All trained models produce scored panel outputs containing predicted probabilities for each pixel--year in the evaluation period. These outputs are used to generate summary metrics, calibration diagnostics, and spatial risk maps. The full pipeline can be re-executed end-to-end given access to the same input datasets and computing environment, enabling exact replication of reported results.

\subsection{Forward predictions}
\label{forward}
To generate probabilistic protection maps for 2025--2030, we distinguish between two model roles. The evaluation model is trained on 2001--2016 and tested on 2017--2019; it produces reported performance metrics and fits the post-hoc probability calibrators (Platt scaling and isotonic regression, ensemble-averaged). The deployment model is then retrained on the full labeled window 2001--2019 using the hyperparameters fixed by the evaluation stage, following standard practice of retraining on all available labeled data once model selection is complete.

The 2019 training cutoff is determined by right-censoring: labeling a pixel as ``protected within 5 years of year $T$'' requires outcome data through $T+5$. Given that WDPA raster data extend to 2024, the last year with a fully observed 5-year outcome window is 2019. Years 2020--2023 are right-censored and excluded from training.

Forward inference is performed on 2024 features. The candidate set consists of all pixels unprotected as of 2023 (WDPA$_{\text{prev}} = 0$). The deployment model assigns each candidate pixel a calibrated probability of becoming protected between 2025 and 2030. This constitutes a strict temporal out-of-sample forecast: no information from 2024 enters model training or calibration.

Three policy scenarios are derived from the resulting probability surface by varying the area threshold applied to top-ranked pixels. The business-as-usual (BAU) scenario selects pixels whose cumulative area equals the observed designation volume between 2019 and 2024 (WDPA$_{2024} -$ WDPA$_{2019}$ in km$^2$), projecting the recent empirical designation rate forward. The moderate scenario targets the midpoint between current coverage and 30\%, and the 30$\times$30 scenario targets the full Kunming--Montreal commitment of 30\% terrestrial protection by 2030.

Methodological validity is established via backtesting: pseudo-forecasts are generated at historical origins $T \in \{2013, 2015, 2017, 2019\}$, each time training on 2001--$(T-1)$, scoring unprotected pixels at year $T$, and evaluating against realized designations in years $T+1$ through $T+5$. Origins $T = 2013$ and $T = 2015$ provide fully observed 5-year outcome windows. This walk-forward evaluation---analogous to expanding-window cross-validation in the time-series literature---provides empirical evidence that the model generalises to future designation events across multiple held-out periods.


\section{Results}
\label{results}

\subsection{Descriptive Statistics and Transition Dynamics}
\label{results:descriptives}

The analysis is based on a spatiotemporal pixel-level panel covering the terrestrial extent of South America from 2000 to 2024 at a $1 \times 1$ km resolution. The final merged dataset contains approximately 520 million pixel-year observations across 20.8 million unique grid cells (equivalent to the square kilometre extent of South America), spanning 25 annual periods. The panel integrates a large set of predictors, described in section~\ref{predictor-variables}.

Data completeness is high across the panel. Missing values are limited and primarily concentrated in ocean-adjacent pixels. Static environmental variables exhibit missingness in roughly 0.5\% of observations, while dynamic predictors such as NDVI, population density, deforestation, and wildfire indicators are missing in fewer than 0.6\% of cases. Spatially smoothed features display even lower missingness rates, reflecting the aggregation across neighboring pixels. The Colombia reference dataset, comprising approximately 24.9 million pixel-year observations across 1.16 million unique pixels, exhibits comparable data quality characteristics.

Protected-area coverage expands substantially over the study period. Across South America, approximately 25.7\% of the terrestrial landscape is designated as protected by 2024. As a consequence of this cumulative expansion, the number of eligible pixels in the risk set declines from roughly 20.8 million in 2000 to about 15.7 million by the end of the sample.

Transitions from unprotected to protected status are defined using a five-year lookahead indicator. Despite the long prediction window, transitions remain rare. Across the full South American panel, the average transition rate is approximately 0.5--0.8\% per pixel-year. Transition probabilities vary substantially over time. During the early part of the sample (2000--2013), transition rates are higher, while later periods exhibit a pronounced decline. In the held-out test period (2017--2019), 380{,}064 transitions occur among roughly 47.7 million pixel-year observations, corresponding to a transition rate below 1\%.

This temporal pattern results in severe and increasing class imbalance. The ratio of non-transition to transition observations rises from roughly 1:30 in the early training period to more than 1:120 in the test period. The Colombia subsample displays distinct dynamics. While transitions are relatively frequent during the early training years, with a substantially more balanced class distribution, transition rates drop sharply after 2014, leading to imbalance levels comparable to those observed at the continental scale.

Taken together, these descriptive patterns highlight three defining features of the data: a very large risk set, strong temporal variation in transition activity, and extreme class imbalance, particularly in recent years. These characteristics shape the empirical challenges addressed in the subsequent predictive analysis.


\subsection{Predictive Performance}
\label{results:performance}

This section evaluates predictive performance on held-out temporal test sets (2017--2019) for both continental (Model~1) and country-level (Model~C) pipelines.

\begin{table}[htbp]
\centering
\caption{Predictive performance across models and spatial scales (temporal test sets only)}
\label{tab:predictive_performance}
\begin{adjustbox}{max width=\textwidth}
\begin{tabular}{lcccccc}
\hline
\textbf{Metric} 
& \textbf{Model 1 (LGBM)} 
& \textbf{Model 1 (RF)} 
& \textbf{Model C (LGBM)} 
& \textbf{Model C (RF)} \\
\hline
ROC--AUC            & 0.9257 & 0.9466 & 0.9464 & 0.9308 \\
PR--AUC             & 0.4299 & 0.4138 & 0.4701 & 0.4353 \\
Precision @ 1\%     & 0.4943 & 0.4613 & 0.5466 & 0.4667 \\
Precision @ 5\%     & 0.1209 & 0.1091 & 0.1578 & 0.1584 \\
Precision @ 10\%    & 0.0640 & 0.0551 & 0.0812 & 0.0807 \\
Lift @ 1\%          & 62.04  & 76.86  & 59.15  & 50.51  \\
Lift @ 5\%          & 15.17  & 18.18  & 17.07  & 17.14  \\
Lift @ 10\%         & 8.03   & 9.18   & 8.79   & 8.73  \\
\hline
\end{tabular}
\end{adjustbox}
\end{table}

Table~\ref{tab:predictive_performance} reports ROC--AUC, PR--AUC, and ranking-based precision metrics across models and spatial scales. Predictive performance is evaluated exclusively on held-out temporal test sets to assess the models’ ability to forecast future protected-area designations under realistic policy conditions. All reported metrics are computed on observations from 2017--2019, with labels capped to avoid right-censoring.

At the continental scale (Model~1), the LightGBM transition model achieves strong discriminative performance, with a ROC--AUC of 0.926 and a PR--AUC of 0.430 on the test set. Given a baseline transition rate of approximately 0.8\%, these values indicate substantial improvement over random ranking. Focusing on policy-relevant thresholds, the model attains a precision of 49.4\% at the top 1\% of predicted risk, corresponding to a lift of roughly 62$\times$ relative to random selection. Precision declines at broader thresholds, reaching 12.1\% at the top 5\% and 6.4\% at the top 10\%, but remains well above baseline rates.

The Random Forest variant of Model~1 exhibits comparable ranking performance. While its ROC--AUC is lower (0.834), its PR--AUC (0.434) and precision at the top 1\% (50.2\%) are similar to those of the LightGBM model, reflecting strong enrichment among the highest-ranked pixels despite reduced overall separability. At wider thresholds, performance patterns are consistent with those of the LightGBM model.

Results are qualitatively similar in the Colombia-specific pipeline (Model~C), albeit at a smaller spatial scale and with different baseline transition dynamics. On the Colombia test set (2017--2019), the LightGBM model achieves a ROC--AUC of 0.946 and a PR--AUC of 0.470, with precision at the top 1\% reaching 54.7\%. The corresponding Random Forest model attains a ROC--AUC of 0.931, a PR--AUC of 0.435, and a top-1\% precision of 46.7\%. These results indicate that both model classes are able to identify a concentrated subset of locations capturing a large share of future designations.

Across both spatial scales and model families, performance is strongest at very high rank thresholds, consistent with the intended use of the models as prioritization tools rather than binary classifiers. The combination of high ROC--AUC, robust PR--AUC, and substantial lift at the top of the prediction distribution suggests that the models successfully exploit spatial, environmental, and socio-economic signals to rank locations by near-term protection likelihood.

\subsection{Aggregation Effects and Event-Level Coverage}
\label{results:aggregation}

Model performance metrics reported throughout this chapter are computed at the pixel--year level, reflecting the underlying panel structure of the data and the forward-looking nature of the prediction task. Under this evaluation scheme, each spatial location may appear multiple times across the test period, and performance metrics such as PR--AUC and Precision@Top-$k$\% quantify the model’s ability to rank pixel--year observations correctly under extreme class imbalance.

However, conservation outcomes are ultimately realized at the \emph{event} or \emph{pixel} level: a location either becomes protected at some point during the test window or it does not. When predictions are aggregated spatially across all test years (2017--2019), a complementary picture emerges. Specifically, a pixel is considered a successful prediction if it is ranked within the top 1\% of predicted risk in \emph{any} test year prior to its designation.

Under this event-level aggregation, Model~1 captures the large majority of protected area establishments. The LightGBM model managed to capture 148{,}673 out of 159{,}350 pixels flagged as high risk subsequently transitioned to protected status, corresponding to a hit rate of 93.3\%. Measured relative to all observed test-period establishments, this implies an event-level recall of 90.4\%. When test-period and post-test establishments (2020--2024) are combined, the model captures approximately two-thirds of all observed designations within the top 1\% of predicted risk. The same pattern is observed for the Random Forest implementation. The RF model managed to capture 152{,}153 out of 159{,}350 pixels flagged as high risk subsequently transitioned to protected status, corresponding to a hit rate of 95.5\%. 


These results clarify the apparent discrepancy between modest row-level PR--AUC values and strong spatial targeting performance. Because protected area designations are temporally clustered and spatially persistent, pixel--year evaluation penalizes models that correctly identify \emph{where} protection will occur but not the exact year of designation. In contrast, pixel-level aggregation reflects the operational question most relevant for conservation planning: whether a location is identified as high risk prior to designation at any point within the planning horizon.

%\subsection{Probability Calibration and Reliability}
%\label{results:calibration}
% Calibration curves, Brier score, before/after calibration comparison.
% Brief and factual.

%\subsection{Spatial Generalisation and Regional Performance}
%\label{results:spatial_generalisation}
% Macro-region results, stability across space.
% Demonstrates model is not driven by one region.

\subsection{Spatial Risk Maps and Predicted Hotspots}
\label{results:spatial_outputs}

Model predictions are translated into spatially explicit risk maps to identify geographic hotspots of future protected-area establishment. For each model, predicted probabilities are ranked and visualised by highlighting the top 1\%, 5\%, and 10\% of pixels across South America. These maps are overlaid with observed protected-area designations from the temporal test period (2017--2019) to facilitate visual comparison between predicted high-risk locations and realised transitions.

At the most restrictive threshold (top 1\%), the risk maps reveal a small number of highly concentrated hotspot regions rather than a diffuse spatial pattern. High-risk pixels tend to form contiguous clusters, reflecting both the spatial structure of the underlying predictors and the clustering of observed protection events. This concentration is consistent across model classes and is particularly pronounced in the LightGBM implementation of Model~1.

Expanding the threshold to the top 5\% and 10\% of predictions increases spatial coverage while preserving clear spatial structure. At these broader thresholds, predicted hotspots extend outward from the highest-risk cores, capturing a larger share of observed designations at the cost of reduced precision. The resulting maps illustrate the trade-off between spatial coverage and targeting intensity that underlies threshold selection in applied conservation settings.

\begin{figure}[H]
\centering
\includegraphics[
  width=\textwidth,
  height=0.8\textheight,
  keepaspectratio
]{model1_lgbm_risk_map_top1pct.pdf}
\caption{Spatial risk map for Model~1 (LightGBM), highlighting the top 1\% of predicted pixels across South America. Shaded areas indicate high-risk locations, overlaid with realised protected-area designations during the temporal test period (2017--2019).}
\label{fig:riskmap_top1}
\end{figure}

\begin{figure}[H]
\centering
\includegraphics[
  width=\textwidth,
  height=0.8\textheight,
  keepaspectratio
]{model1_rf_risk_map_top1pct.pdf}
\caption{Spatial risk map for Model~1 (Random Forest), highlighting the top 1\% of predicted pixels across South America. Shaded areas indicate high-risk locations, overlaid with realised protected-area designations during the temporal test period (2017--2019).}
\label{fig:riskmap_top1_rf}
\end{figure}

\begin{figure}[H]
\centering
\includegraphics[
  width=\textwidth,
  height=0.8\textheight,
  keepaspectratio
]{model1_lgbm_risk_map_top5pct.pdf}
\caption{Spatial risk map for Model~1 (LightGBM), highlighting the top 5\% of predicted pixels across South America. Shaded areas indicate high-risk locations, overlaid with realised protected-area designations during the temporal test period (2017--2019).}
\label{fig:riskmap_top5_model1_lgbm}
\end{figure}

\begin{figure}[H]
\centering
\includegraphics[
  width=\textwidth,
  height=0.8\textheight,
  keepaspectratio
]{model1_lgbm_risk_map_top10pct.pdf}
\caption{Spatial risk map for Model~1 (LightGBM), highlighting the top 10\% of predicted pixels across South America. Shaded areas indicate high-risk locations, overlaid with realised protected-area designations during the temporal test period (2017--2019).}
\label{fig:riskmap_top10_model1_lgbm}
\end{figure}

\begin{figure}[H]
\centering
\includegraphics[
  width=\textwidth,
  height=0.8\textheight,
  keepaspectratio
]{modelc_rf_risk_map_top1pct.pdf}
\caption{Spatial risk map for Model~C (Colombia; Random Forest), highlighting the top 1\% of predicted pixels. Shaded areas indicate high-risk locations, overlaid with realised protected-area designations during the temporal test period (2017--2019).}
\label{fig:riskmap_top1_modelc_lgbm}
\end{figure}

Probability maps displaying continuous predicted transition probabilities across the full landscape show a strongly right-skewed distribution. The vast majority of pixels are assigned very low transition probabilities, while a small upper tail concentrates substantially higher relative risk. High-probability pixels are spatially autocorrelated, forming coherent patches rather than isolated points, which reinforces the interpretation of predictions as relative spatial rankings rather than independent point estimates.

\begin{figure}[H]
\centering
\includegraphics[
  width=\textwidth,
  height=0.8\textheight,
  keepaspectratio
]{model1_lgbm_probability_map.pdf}
\caption{Continuous predicted transition probabilities from Model~1 (LightGBM) across South America. The distribution is strongly right-skewed: most pixels receive near-zero probability, while a small upper tail concentrates substantially higher relative risk.}
\label{fig:probmap_model1_lgbm}
\end{figure}

\begin{figure}[H]
\centering
\includegraphics[
  width=\textwidth,
  height=0.8\textheight,
  keepaspectratio
]{modelc_lgbm_probability_map.pdf}
\caption{Continuous predicted transition probabilities from Model~C (Colombia; LightGBM). High-probability pixels cluster spatially, forming coherent patches rather than isolated points, consistent with the interpretation of model outputs as relative spatial rankings.}
\label{fig:probmap_modelc_lgbm}
\end{figure}

Overall, the spatial outputs demonstrate that the models generate interpretable and spatially coherent risk surfaces. By translating statistical predictions into geographically explicit maps, the framework enables direct visual identification of priority areas and provides a transparent link between model outputs and on-the-ground conservation planning.

\subsection{Key Predictors and Model Drivers}
\label{results:drivers}

To interpret what drives predicted \emph{new protected-area designations} over a five-year lookahead, this section summarizes feature importance patterns from the two final modelling pipelines: Model~1 (continent-scale South America) and Model~C (Colombia). Feature-importance values are not directly comparable across algorithms. Random Forest (RF) relies on impurity-based importance, while LightGBM (LGBM) reports tree-based importance from the fitted gradient-boosted model. We therefore focus on robust qualitative regularities: which predictors consistently appear among the most influential, and how emphasis differs across spatial scope.

\subsubsection*{Accessibility, infrastructure, and human presence}

Across both spatial scopes, predictors that proxy \emph{accessibility and human activity} are central. In the continent-scale RF model, \texttt{dist\_road} is the single most influential feature, indicating that proximity to road networks is strongly informative for where new designations occur. In the corresponding LGBM model, multiple infrastructure-distance variables rank among the top drivers, including \texttt{dist\_powerplant} and \texttt{dist\_oil\_gas}. In parallel, population density (\texttt{GPW\_b1} and smoothed variants) and night-time lights (\texttt{HNTL\_b1\_smooth64}) repeatedly appear among top predictors, consistent with designation processes being tightly linked to human settlement patterns and land-use pressure.

\subsubsection*{Landscape dynamics and path dependence}

The strongest and most stable predictor across both LGBM models is past deforestation intensity aggregated at broader spatial scales: \texttt{deforestation\_b1\_smooth64} ranks first in Model~1 LGBM and Model~C LGBM. This suggests strong spatial clustering in designation events: areas embedded in, or adjacent to, zones of recent land conversion are more likely to be designated within the subsequent five-year window. Vegetation productivity signals (\texttt{NDVI\_b1\_smooth64}) also rank highly in both LGBM runs, implying that canopy/greenness conditions contribute to where designations occur, conditional on anthropogenic context.

\subsubsection*{Protected-area network proximity and regional differences}

Proximity to the existing protected-area network (\texttt{dist\_wdpa}) is a consistently important predictor in both spatial scopes, with particularly strong prominence in Model~C LGBM (ranked second). This pattern is consistent with incremental expansion dynamics: new designations tend to occur near the current network (e.g., boundary extensions, buffers, or corridor-type additions). A notable contrast is that the Colombia RF model places comparatively more emphasis on biophysical and landscape-structure predictors, especially smoothed GSN layers (e.g., \texttt{GSN\_b2\_smooth16/64}) and multiple WorldClim variables, whereas the continent RF model highlights road proximity and land cover alongside population measures. This indicates that within-country heterogeneity may amplify the role of ecological suitability proxies, while continent-wide patterns are more strongly dominated by accessibility gradients.

\begin{table}[H]
\centering
\small
\caption{Top 10 predictors by model-reported importance across continent-scale (Model~1) and Colombia-specific (Model~C) models. Rankings are shown separately for Random Forest (RF) and LightGBM (LGBM). Importance values are not directly comparable across algorithms.}
\label{tab:top_features}
\resizebox{\textwidth}{!}{
\begin{tabular}{r|ll|ll}
\hline
\textbf{Rank} 
& \multicolumn{2}{c|}{\textbf{Model 1 (South America)}} 
& \multicolumn{2}{c}{\textbf{Model C (Colombia)}} \\
& \textbf{RF} & \textbf{LGBM} & \textbf{RF} & \textbf{LGBM} \\
\hline
1  & dist\_road                  & deforestation\_b1\_smooth64 & GSN\_b2\_smooth16        & deforestation\_b1\_smooth64 \\
2  & GPW\_b1                     & dist\_wdpa                  & GSN\_b2\_smooth64        & dist\_wdpa                  \\
3  & landcover\_b1               & dist\_oil\_gas               & GSN\_b2                  & dist\_oil\_gas               \\
4  & GPW\_b1\_smooth16            & GPW\_b1                     & WorldClim\_b17           & GPW\_b1                     \\
5  & GSN\_b2\_smooth16            & NDVI\_b1\_smooth64           & dist\_powerplant         & NDVI\_b1\_smooth64           \\
6  & elevation\_b2               & dist\_powerplant             & HNTL\_b1\_smooth64       & HNTL\_b1\_smooth64           \\
7  & NDVI\_b1\_smooth64           & WorldClim\_b14               & WorldClim\_b11           & WorldClim\_b17               \\
8  & dist\_wdpa                  & GSN\_b2\_smooth64            & GPW\_b1\_smooth16        & WorldClim\_b15               \\
9  & HNTL\_b1\_smooth64           & wildfire\_b1\_smooth64       & WorldClim\_b14           & WorldClim\_b14               \\
10 & WorldClim\_b12               & GPW\_b1\_smooth64            & dist\_wdpa               & WorldClim\_b7                \\
\hline
\end{tabular}
}
\end{table}

\subsubsection*{Environmental constraints}

Climatic predictors from WorldClim (e.g., \texttt{WorldClim\_b14}, \texttt{WorldClim\_b17}, \texttt{WorldClim\_b11}, \texttt{WorldClim\_b7}) and topography (e.g., \texttt{elevation\_b2} in Model~C LGBM) frequently appear among top-ranked features. These variables likely provide a secondary constraint set: conditional on accessibility and land-use context, environmental suitability helps discriminate which areas are more likely to be designated.

\subsubsection*{Synthesis}

Taken together, the driver patterns suggest a consistent hierarchy. First, designation likelihood is strongly shaped by \emph{human accessibility and pressure} (roads, population, night-time lights, and infrastructure distances). Second, it exhibits \emph{spatial path dependence} through strong associations with recent deforestation dynamics and multi-scale neighbourhood context. Third, \emph{network proximity} (distance to existing protected areas) indicates that new designations are frequently spatially adjacent to the current protected-area system. Finally, climate and topography act as secondary constraints that refine where designations occur within accessible landscapes.

\subsection{Spatial CV}
\label{spatialcv}

\noindent \textbf{Layer 2 -- Biome-stratified evaluation.} Using the full production models scored on the 2017--2019 test period, we evaluated performance separately for each WWF terrestrial biome via GSN ecoregion lookup.

\medskip

\noindent For South America (Model 1; 47.7M pixel-years, 380,064 positives), both models perform well in aggregate (LGBM: ROC-AUC $= 0.875$, PR-AUC $= 0.634$; RF: ROC-AUC $= 0.927$, PR-AUC $= 0.634$), with RF showing consistently higher ROC-AUC across most biomes. Performance is strongest in Deserts \& Xeric Shrublands (RF ROC-AUC $= 0.999$), Mangroves (0.993), and Tropical \& Subtropical Dry Broadleaf Forests (0.980). The weakest biome is Montane Grasslands \& Shrublands (Andean highlands; LGBM 0.629, RF 0.590), reflecting an extremely sparse positive rate (0.07\%) and likely heterogeneous, politically complex protection dynamics.

\medskip

\noindent For the USA (Model 2; 36.8M pixel-years, 14,787 positives), the RF substantially outperforms LGBM in terms of PR-AUC (0.559 vs. 0.082), despite comparable ROC-AUC (0.999 vs. 0.944). This divergence likely reflects the extreme class imbalance in the USA test set (0.04\% positive rate), under which LGBM's probability calibration struggles to rank the small number of positives highly. 85\% of all USA positives fall within Deserts \& Xeric Shrublands (BLM land in the American West), making the evaluation heavily dominated by a single biome type. Four biomes --- Flooded Grasslands \& Savannas, Mangroves, Tropical \& Subtropical Coniferous Forests, and Tropical \& Subtropical Grasslands, Savannas \& Shrublands --- recorded zero new protected area designations in the 2017--2019 test period and could not be evaluated.

\medskip

\noindent \textbf{Layer 3 -- Cross-continental transfer.} We applied each production model to the opposite continent without retraining (SA $\rightarrow$ USA; USA $\rightarrow$ SA), using the identical 73-feature set. Transfer performance is near-random across all four combinations (ROC-AUC range: 0.301--0.531; PR-AUC $< 0.01$ in all cases). Notably, the SA-trained RF applied to the USA is anti-predictive (ROC-AUC $= 0.301$), assigning systematically higher scores to pixels that did not become protected. These results indicate that protection drivers are strongly region-specific: the models capture institutional, political, and land-use patterns that do not generalise across continental governance contexts. This finding is itself substantively informative --- it implies that the $30 \times 30$ expansion pathway is shaped by continent-level policy environments rather than universal ecological or geographical characteristics.

\vspace{1em}

Layer 2 tells you the model is discriminative within the geography it was trained on. RF AUC 0.927 (SA) and 0.999 (USA) overall are strong. The model has genuinely learned to rank pixels. But this is the easy case --- training and test data are from the same continent, same feature distribution.

\vspace{1em}

Layer 1 (LOBO) is the real generalization test. The gap between Layer 2 and Layer 1 is dramatic, and tells a clear story:

\begin{itemize}
    \item \textbf{Deserts \& Xeric Shrublands in SA (LOBO AUC 0.996)} is a strong positive result. The model generalises perfectly to this biome from training on all others. Why? Because this biome is almost entirely the Atacama Desert --- PA designation there follows very distinctive patterns (mineral zones, extreme remoteness) that other biome features predict well. Also, this biome has 109K positives --- by far the richest test set.
    
    \item \textbf{Temperate Broadleaf in SA (LOBO AUC 0.460)} is sub-random. The model trained on the other 8 SA biomes cannot predict PA designation in Patagonian temperate forests. This biome follows political/conservation patterns orthogonal to the rest of SA.
    
    \item \textbf{USA LOBO at mean AUC 0.54} is near-random across all biomes. The USA model learns from its geography as a whole, but no individual biome can be reliably predicted when held out. This is partly a class imbalance issue (only 329K positives spread across 12 biomes, with 85\% in deserts), but also genuine: USA PA designation is driven by jurisdiction and political factors that look very different across biomes.
\end{itemize}

\vspace{1em}

Layer 3 (cross-continental) fails completely. AUC $\sim$0.51--0.53 across all four directions. This is your strongest methodological finding: the patterns the models learn are continent-specific, not universal. Protected area designation in the Amazon follows entirely different rules than designation in Montana. Your features (which are generic --- elevation, climate, biodiversity, infrastructure distance) are insufficient to capture the institution-specific, jurisdiction-specific, and political-economy-specific drivers of PA designation. The model is not learning ``what makes a pixel become protected'' in a general sense --- it's learning ``what kinds of land have historically been designated in THIS region.''

\section{Discussion and Conclusion}
\label{discussion}

This thesis examined whether future protected-area (PA) establishments can be forecast using high-resolution spatial data and machine-learning models, while accounting for the rarity and policy-driven nature of conservation outcomes. By framing PA designation as a transition within a clearly defined risk set and evaluating performance exclusively on temporally held-out data, the analysis shows that future PA placements are predictable to a meaningful extent at both continental and national scales.

Across South America, LightGBM and Random Forest models achieved strong temporal out-of-sample performance, particularly in ranking pixels by relative transition risk. A small fraction of the landscape, most notably the top 1\% of predictions, captured a large share of subsequent PA establishments, indicating substantial concentration of predictive signal. Spatial risk maps revealed coherent hotspot regions rather than diffuse patterns, supporting the interpretation of model outputs as structured spatial rankings rather than noise. Comparable results obtained in the Colombia-specific pipeline suggest that these findings are not driven solely by continental aggregation.

In conclusion, this thesis demonstrates that machine-learning methods can generate informative, spatially explicit forecasts of future PA expansion when applied within a rigorous evaluation framework. By combining high-resolution data, risk-set logic, and ranking-based metrics, the approach offers a transparent decision-support tool, while highlighting the importance of cautious interpretation when extrapolating across space or policy regimes.

\section{References}\label{references}
Baldi, G., Texeira, M., Martin, O. A., Grau, H. R., \& Jobbágy, E. G. (2017). Opportunities drive the global distribution of protected areas. PeerJ 5, e2989.
\newline \newline
Barnes, M. Protect biodiversity, not just area. Nature 526, 195 (2015). https://doi.org/10.1038/526195e
\newline \newline
Bühler, C. K., \& Benson, H. Y. (2024, March). Decision-making for land conservation: A derivative-free optimization framework with nonlinear inputs. Proceedings of the AAAI Conference on Artificial Intelligence, 38(20), 21932–21939. https://doi.org/10.1609/aaai.v38i20.30195
\newline \newline
Beckmann, M., Gerstner, K., Akin-Fajiye, M., Ceaușu, S., Kambach, S., Kinlock, N. L., Phillips, H. R. P., Verhagen, W., Gurevitch, J., Klotz, S., Newbold, T., Verburg, P. H., Winter, M., \& Seppelt, R. (2019). Conventional land-use intensification reduces species richness and increases production: A global meta-analysis. Global Change Biology, 25(6), 1941–1956. https://doi.org/10.1111/gcb.14606
\newline \newline
Berglund, Moa, Martin Nilsson Jacobi, and Per R. Jonsson. "Optimal selection of marine protected areas based on connectivity and habitat quality." Ecological Modelling 240 (2012): 105-112.
\newline \newline
Branquart, E., Verheyen, K., \& Latham, J. (2008). Selection criteria of protected forest areas in Europe: the theory and the real world. Biological conservation, 141(11), 2795-2806.
\newline \newline
Brondizio, E., Díaz, S. M., Settele, J., Ngo, H., Gueze, M., Aumeeruddy-Thomas, Y., ... \& Zayas, C. (2019). Assessing a planet in transformation: Rationale and approach of the IPBES Global Assessment on Biodiversity and Ecosystem Services.
\newline \newline
Brown, C. F., Kazmierski, M. R., Pasquarella, V. J., Rucklidge, W. J., Samsikova, M., Zhang, C., ... \& Kohli, P. (2025). AlphaEarth Foundations: An embedding field model for accurate and efficient global mapping from sparse label data. arXiv preprint arXiv:2507.22291.
\newline \newline
Chape S, Harrison J, Spalding M, Lysenko I. Measuring the extent and effectiveness of protected areas as an indicator for meeting global biodiversity targets. Philos Trans R Soc Lond B Biol Sci. 2005 Feb 28;360(1454):443-55. doi: 10.1098/rstb.2004.1592. PMID: 15814356; PMCID: PMC1569446.
\newline \newline
Chape, Stuart, Jenkins, Martin, \& Spalding, Mark. (2008). The world’s protected areas : status, values and prospects in the 21st Century. UNEP-WCMC. https://www.biodiversitylibrary.org/item/119102
\newline \newline
Chaplin-Kramer, R., Neugarten, R.A., Sharp, R.P. et al. Mapping the planet’s critical natural assets. Nat Ecol Evol 7, 51–61 (2023). https://doi.org/10.1038/s41559-022-01934-5
\newline \newline
Convention on Biological Diversity. (2022). Kunming‐Montreal global biodiversity framework. In Convention on Biological Diversity. Montreal, Canada: Convention on Biological Diversity.
\newline \newline
Dasgupta, S., Blankespoor, B., \& Wheeler, D. (2025). Implementing 30x30: Lessons from country case studies (Policy Research Working Paper No. 11045). World Bank. http://hdl.handle.net/10986/42725
\newline \newline
Doyle, C., Beach, T., \& Luzzadder-Beach, S. (2021). Tropical Forest and Wetland Losses and the Role of Protected Areas in Northwestern Belize, Revealed from Landsat and Machine Learning. Remote Sensing, 13(3), 379. https://doi.org/10.3390/rs13030379
\newline \newline
Nigel Dudley and Sue Stolton (eds) (2008). Defining protected areas: an international conference in Almeria, Spain. Gland, Switzerland: IUCN. 220 pp 
\newline \newline
da Silveira, C. B. L., Strenzel, G. M. R., Maida, M., Gaspar, A. L. B., \& Ferreira, B. P. (2021). Coral Reef Mapping with Remote Sensing and Machine Learning: A Nurture and Nature Analysis in Marine Protected Areas. Remote Sensing, 13(15), 2907. https://doi.org/10.3390/rs13152907
\newline \newline
Dudley, N., and Stolton, S. (eds.). 2022. Best Practice in Delivering the 30x30 Target (2nd Edition, October 2022). The Nature Conservancy and Equilibrium Research.
\newline \newline
Sandra Díaz et al. ,Pervasive human-driven decline of life on Earth points to the need for transformative change.Science366,eaax3100(2019).DOI:10.1126/science.aax3100
\newline \newline
Eken, G., L. Bennun, and C. Boyd. 2004b. Protected areas design and systems planning: Key requirements for successful planning, site selection and establishment of protected areas. Pages. 37–44. In the Secretariat of the Convention on Biological Diversity (SCBD). Biodiversity Issues Consideration in the Planning, Establishment and Management of Protected Area Sites and Networks. Montreal: SCBD.
\newline \newline
Fahrig, L. (2003). Effects of habitat fragmentation on biodiversity. Annual review of ecology, evolution, and systematics, 34(1), 487-515.
\newline \newline
Fisher, J., Allen, S., Yetman, G. et al. Assessing the influence of landscape conservation and protected areas on social wellbeing using random forest machine learning. Sci Rep 14, 11357 (2024). https://doi.org/10.1038/s41598-024-61924-4
\newline \newline
Fortin, D., et al. (2020). When the protection of a threatened species depends on the economy of a foreign nation. PLOS ONE, 15(3), e0229555. https://doi.org/10.1371/journal.pone.0229555
\newline \newline
Gaston, Kevin J., et al. "The ecological performance of protected areas." Annual review of ecology, evolution, and systematics 39.1 (2008): 93-113.
\newline \newline
Gaston, K.J., Jackson, S.F., Nagy, A., Cantú-Salazar, L. and Johnson, M. (2008), Protected Areas in Europe. Annals of the New York Academy of Sciences, 1134: 97-119. https://doi.org/10.1196/annals.1439.006
\newline \newline
Gibson, J., Olivia, S., Boe-Gibson, G., \& Li, C. (2021). Which night lights data should we use in economics, and where?. Journal of Development Economics, 149, 102602.
\newline \newline
Girardin, C. A., Jenkins, S., Seddon, N., Allen, M., Lewis, S. L., Wheeler, C. E., ... \& Malhi, Y. (2021). Nature-based solutions can help cool the planet—if we act now. Nature, 593(7858), 191-194.
\newline \newline
Govind, P. J., \& Webb, T. (2025). The nature of nature: Why the promotion of plural values of nature could pose transition risks and influence nature‐related financial disclosure. Review of European, Comparative \& International Environmental Law.
\newline \newline
Greenhill, S., et al. (2024). Machine learning predicts which rivers, streams, and wetlands the Clean Water Act regulates. Science, 383(6674), 406–412. https://doi.org/10.1126/science.adi3794
\newline \newline
He, K., Zhang, X., Ren, S., \& Sun, J. (2016). Deep residual learning for image recognition. In 2016 IEEE Conference on Computer Vision and Pattern Recognition (CVPR) (pp. 770–778). https://doi.org/10.1109/CVPR.2016.90
\newline \newline
Hooper, D.U., Chapin, F.S., Ewel, J.J., Hector, A., Inchausti, P., Lavorel, S., Lawton, J.H., Lodge, D.M., Loreau, M., Naeem, S., Schmid, B., Set¨ ala, ¨ H., Symstad, A.J., Vandermeer, J., Wardle, D.A., 2005. Effects of biodiversity on ecosystem functioning: a consensus of current knowledge. Ecol. Monogr. 75, 3–35. https://doi. org/10.1890/04-0922.
\newline \newline
Hughes, T., Kerry, J., Álvarez-Noriega, M. et al. Global warming and recurrent mass bleaching of corals. Nature 543, 373–377 (2017). https://doi.org/10.1038/nature21707
\newline \newline
Jinhoo Hwang, Yuyoung Choi, Yoonji Kim, Lim No Ol, Young-Jae Yoo, Hyo Jin Cho, Zhemin Sun, Seongwoo Jeon, Analysis of the effect of environmental protected areas on land-use and carbon storage in a megalopolis, Ecological Indicators, Volume 133, 2021, 108352, ISSN 1470-160X, https://doi.org/10.1016/j.ecolind.2021.108352.
\newline \newline
International Union for Conservation of Nature. World Conservation Strategy: Living Resource Conservation for Sustainable Development (IUCN–UNEP–WWF, 1980). 
\newline \newline
Jeremias, G., Barbosa, J., Marques, S.M., Asselman, J., Goncalves, F.J.M., Pereira, J.L., 2018. Synthesizing the role of epigenetics in the response and adaptation of species to climate change in freshwater ecosystems. Mol. Ecol. 27 (13), 2790–2806. https:// doi.org/10.1111/mec.14727. 
\newline \newline
Joppa LN, Pfaff A (2009) High and Far: Biases in the Location of Protected Areas. PLoS ONE 4(12): e8273. doi:10.1371/journal.pone.0008273
\newline \newline
Juffe-Bignoli, D., Burgess, N.D., Bingham, H., Belle, E.M.S., de Lima, M.G., Deguignet, M., Bertzky, B., Milam, A.N., Martinez-Lopez, J., Lewis, E., Eassom, A., Wicander, S., Geldmann, J., van Soesbergen, A., Arnell, A.P., O’Connor, B., Park, S., Shi, Y.N., Danks, F.S., MacSharry, B., Kingston, N. (2014). Protected Planet Report 2014. UNEP-WCMC: Cambridge, UK. 
\newline \newline
Kleinberg, J., Lakkaraju, H., Leskovec, J., Ludwig, J., \& Mullainathan, S. (2018). Human decisions and machine predictions. The Quarterly Journal of Economics, 133(1), 237–293. https://doi.org/10.1093/qje/qjx032
\newline \newline
Lago-González, P., Acuña Alonso, C. \& Álvarez, X. Influence of Protected Areas in the Carbon Cycle: A Machine Learning Approach for Future Forecasting. Earth Syst Environ 8, 1683–1697 (2024). https://doi.org/10.1007/s41748-024-00451-8
\newline \newline
Ma, B., Zhao, Z., Ding, H. et al. Household costs and benefits of biodiversity conservation: case study of Sichuan giant panda reserves in China. Environ Dev Sustain 20, 1665–1686 (2018). https://doi.org/10.1007/s10668-017-9959-z
\newline \newline
Mayfield, H., Smith, C., Gallagher, M., \& Hockings, M. (2017). Use of freely available datasets and machine learning methods in predicting deforestation. Environmental modelling \& software, 87, 17-28.
\newline \newline
McKinney, M.L. Effects of urbanization on species richness: A review of plants and animals. Urban Ecosyst 11, 161–176 (2008). https://doi.org/10.1007/s11252-007-0045-4
\newline \newline
Melillo, J.M., Lu, X., Kicklighter, D.W. et al. Protected areas’ role in climate-change mitigation. Ambio 45, 133–145 (2016). https://doi.org/10.1007/s13280-015-0693-1
\newline \newline
Miranda, J., Britz, W., \& Börner, J. (2024). Impacts of commodity prices and governance on the expansion of tropical agricultural frontiers. Scientific Reports, 14, 9209. https://doi.org/10.1038/s41598-024-59446-0
\newline \newline
Millennium Ecosystem Assessment, 2005. Ecosystems and Human Well-being: Opportunities and Challenges for Business and Industry. World Resources Institute, Washington, DC.
\newline \newline
Mouillot, D., Velez, L., Albouy, C., \& others. (2024). The socioeconomic and environmental niche of protected areas reveals global conservation gaps and opportunities. Nature Communications, 15, 9007. https://doi.org/10.1038/s41467-024-53241-1
\newline \newline
Mühling, Ş. M. K. (2023). Utilizing artificial intelligence (AI) for the identification and management of marine protected areas (MPAs): a review. Journal of Geoscience and Environment Protection, 11(9), 118-132.
\newline \newline
Ndayizeye, G., Imani, G., Nkengurutse, J., Irampagarikiye, R., Ndihokubwayo, N., Niyongabo, F., \& Cuni-Sanchez, A. (2020). Ecosystem services from mountain forests: Local communities’ views in Kibira National Park, Burundi. Ecosystem Services, 45, 101171.
\newline \newline
Nogueira, E.M., Yanai, A.M., de Vasconcelos, S.S. et al. Carbon stocks and losses to deforestation in protected areas in Brazilian Amazonia. Reg Environ Change 18, 261–270 (2018). https://doi.org/10.1007/s10113-017-1198-1
\newline \newline
Robert J. Nicholls, Anny Cazenave ,Sea-Level Rise and Its Impact on Coastal Zones.Science328,1517-1520(2010).DOI:10.1126/science.1185782
\newline \newline
Ignacio Palomo, Berta Martín-López, Marion Potschin, Roy Haines-Young, Carlos Montes, National Parks, buffer zones and surrounding lands: Mapping ecosystem service flows, Ecosystem Services, Volume 4, 2013, Pages 104-116, ISSN 2212-0416, https://doi.org/10.1016/j.ecoser.2012.09.001.
\newline \newline
Pettorelli, N., Williams, J., Schulte to Buehne, H., \& Crowson, M. (2024, June 17). Deep learning and satellite remote sensing for biodiversity monitoring and conservation. Remote Sensing in Ecology and Conservation, 11, 123–132. https://doi.org/10.1002/rse2.415
\newline \newline
Phillips, A. (2004). The history of protected areas and the IUCN categories system. PARKS: The International Journal for Protected Area Managers, 14(3), 4–14.
\newline \newline
Plas, T. van der, Alexander, D. G., \& Pocock, M. J. O. (2025, May 12). Monitoring protected areas by integrating machine learning, remote sensing and citizen science. Ecological Solutions and Evidence, 6, Article e70040. https://doi.org/10.1002/2688-8319.70040
\newline \newline
Prudhomme, R., De Palma, A., Dumas, P., Gonzalez, R., Leadley, P., Levrel, H., Brunelle, T., 2020. Combining mitigation strategies to increase co-benefits for biodiversity and food security. Environ. Res. Lett. 15 (11) https://doi.org/10.1088/1748-9326/abb10a. 
\newline \newline
Reynaert, M., Souza-Rodrigues, E., \& van Benthem, A. A. (2024). The environmental impacts of protected area policy. Regional Science and Urban Economics, 107, 103968. https://doi.org/10.1016/j.regsciurbeco.2023.103968
\newline \newline
Robinson, J. G., LaBruna, D., O’Brien, T., Clyne, P. J., Dudley, N., Andelman, S. J., ... \& Watson, J. E. (2024). Scaling up area-based conservation to implement the Global Biodiversity Framework’s 30x30 target: The role of Nature’s Strongholds. PLoS Biology, 22(5), e3002613.
\newline \newline
Rolf, E., Proctor, J., Carleton, T., et al. (2021). A generalizable and accessible approach to machine learning with global satellite imagery. Nature Communications, 12, 4392. https://doi.org/10.1038/s41467-021-24638-z
\newline \newline
Ruffault, J., et al. (2018). Predicting habitat quality of protected dry grasslands using Landsat NDVI phenology. Ecological Indicators, 91, 447–460. https://doi.org/10.1016/j.ecolind.2018.04.012
\newline \newline
Salvio, G. M. M., \& Gomes, C. R. (2018). Protected area systems in south American countries. Floresta e Ambiente, 25(4), e20171134.
\newline \newline
Schippers, P., van der Heide, C.M., Koelewijn, H.P. et al. Landscape diversity enhances the resilience of populations, ecosystems and local economy in rural areas. Landscape Ecol 30, 193–202 (2015). https://doi.org/10.1007/s10980-014-0136-6
\newline \newline
Uta Schirpke, Genxu Wang, Emilio Padoa-Schioppa, Editorial: Mountain landscapes: Protected areas, ecosystem services, and future challenges, Ecosystem Services, Volume 49, 2021, 101302, ISSN 2212-0416, https://doi.org/10.1016/j.ecoser.2021.101302.
\newline \newline
Schuster, R., Buxton, R., Hanson, J. O., Binley, A. D., Pittman, J., Tulloch, V., ... \& Bennett, J. R. (2023). Protected area planning to conserve biodiversity in an uncertain future. Conservation Biology, 37(3), e14048.
\newline \newline
Sharma, R., Jones, S., Bastin, L., \& Gordon, A. (2025). Private protected areas exhibit greater bias towards unproductive land compared to public protected areas. Environmental Research Communications, 7, 041003. https://doi.org/10.1088/2515-7620/adc546
\newline \newline
Shin, Y.-J., Midgley, G.F., Archer, E.R.M., Arneth, A., Barnes, D.K.A., Chan, L., Smith, P., 2022. Actions to halt biodiversity loss generally benefit the climate. Glob. Chang. Biol. 28 (9), 2846–2874. https://doi.org/10.1111/gcb.16109. 
\newline \newline
Steif, K. (n.d.). A predictive site suitability approach to land conservation (Pennsylvania example). Urban Spatial Analysis [Blog]. https://urbanspatialanalysis.com/a-predictive-site-suitability-approach-to-land-conservation/
\newline \newline
TDK Corporation. (2024). TNFD disclosure: Dependence and impact on nature (Protected areas as transition risks).
\newline \newline
Valente, J.J., Gannon, D.G., Hightower, J. et al. Toward conciliation in the habitat fragmentation and biodiversity debate. Landsc Ecol 38, 2717–2730 (2023). https://doi.org/10.1007/s10980-023-01708-9
\newline \newline
Van der Plas, T. L., Alexander, D. G., \& Pocock, M. J. O. (2025). Monitoring protected areas by integrating machine learning, remote sensing and citizen science. Ecological Solutions and Evidence, 6(2), Article e70040. https://doi.org/10.1002/2688-8319.70040
\newline \newline
Venter, O., Magrach, A., Outram, N., Klein, C.J., Possingham, H.P., Di Marco, M. and Watson, J.E.M. (2018), Bias in protected-area location and its effects on long-term aspirations of biodiversity conventions. Conservation Biology, 32: 127-134. https://doi.org/10.1111/cobi.12970
\newline \newline
Zhirong Wang, Tongxin Wang, Xiujuan Zhang, Junbang Wang, Yongsheng Yang, Yu Sun, Xiaohua Guo, Qinghua Wu, Eugenie Nepovimova, Alan E. Watson, Kamil Kuca, Biodiversity conservation in the context of climate change: Facing challenges and management strategies, Science of The Total Environment, Volume 937, 2024, 173377, ISSN 0048-9697, https://doi.org/10.1016/j.scitotenv.2024.173377.
\newline \newline
Watson, J., Dudley, N., Segan, D. et al. The performance and potential of protected areas. Nature 515, 67–73 (2014). https://doi.org/10.1038/nature13947
\newline \newline
Wheeler t., \& von Braun, J. (2013). climate change impacts on global food security. Science (New York, N.Y.), 341(6145),508–513. https://doi.org/10.1126/science.1239402
\newline \newline
D.R. Williams,C. Rondinini, \& D. Tilman,  Global protected areas seem insufficient to safeguard half of the world's mammals from human-induced extinction, Proc. Natl. Acad. Sci. U.S.A. 119 (24) e2200118119, https://doi.org/10.1073/pnas.2200118119 (2022).
\newline \newline
UNEP-WCMC and IUCN (2024). Protected Planet Report 2024. UNEP-WCMC and IUCN: Cambridge, United Kingdom; Gland, Switzerland.
\newline \newline
Xu, B., Pan, J. Spatial distribution characteristics of national protected areas in China. J. Geogr. Sci. 29, 2047–2068 (2019). https://doi.org/10.1007/s11442-019-1704-0
\newline \newline
Xu, H., Cao, Y., Yu, D. et al. Ensuring effective implementation of the post-2020 global biodiversity targets. Nat Ecol Evol 5, 411–418 (2021). https://doi.org/10.1038/s41559-020-01375-y
\newline \newline
W. Xu, Y. Xiao, J. Zhang, W. Yang, L. Zhang, V. Hull, Z. Wang, H. Zheng, J. Liu, S. Polasky, L. Jiang, Y. Xiao, X. Shi, E. Rao, F. Lu, X. Wang, G.C. Daily, \& Z. Ouyang, Strengthening protected areas for biodiversity and ecosystem services in China, Proc. Natl. Acad. Sci. U.S.A. 114 (7) 1601-1606, https://doi.org/10.1073/pnas.1620503114 (2017).
\newline \newline
Yang, Hongbo, et al. "Feedback of telecoupling: the case of a payments for ecosystem services program." Ecology \& Society 23.2 (2018).
\newline \newline


\clearpage


\setstretch{1.5}
\setcitestyle{authoryear}
\crefalias{section}{appsec}
\setcounter{table}{0}
\renewcommand\thetable{\Alph{section}.\arabic{table}}
\setcounter{equation}{0}
\renewcommand\theequation{\Alph{section}.\arabic{equation}}
\setcounter{figure}{0}
\renewcommand\thefigure{\Alph{section}.\arabic{figure}}

\begin{appendices}

\section{Declaration on the Use of Artificial Intelligence}
\label{declaration }
In the preparation of this thesis, the artificial intelligence tools of Open-AI Chat-GPT, Cursor, and Overleaf were used to support various aspects of the writing and analysis process. These tools assisted in language refinement, improving clarity of expression, and verifying whether formulations were understandable to a broader audience. They also served as a coding support tool, helping to brainstorm implementation strategies, to improve coding techniques, and to cross-check the logic of scripts and outputs.

All research design, data analysis, modeling, interpretation of results, and final text composition were carried out independently by the author. The thesis reflects the author’s own scholarly work, critical thinking, and academic judgment.
\end{appendices}

\end{document}
